%%

%

%%%%

%%%   Самарский государственный технический университет

%%%   Кафедра прикладной математики и информатики

%%%

%%%   Преамбула для статьи в журнал "Вестник Самарского государственного

%%%   технического университета. Серия: «Физико-математические науки»"

%%%   Ориентирована под MikTEx 2.4 for Win32

%%%

%%%   Хотя сам журнал собирается под GNU/Linux на дистрибутиве TeXLive



\documentclass[11pt,a4paper,twoside]{article}



\usepackage[cp1251]{inputenc}

\usepackage[T2A]{fontenc}

\usepackage[english,russian]{babel}

\usepackage{samgtubib}

\usepackage[dvips]{graphicx}

\usepackage[multiple]{footmisc}



\usepackage{samgtu}

\renewcommand{\VestnikYear}{2009} % Год издания Вестника

\renewcommand{\VestnikNumber}{2\,(19)} % Номер Вестника

\renewcommand{\VestnikMonth}{Ноябрь} % Месяц выхода Вестника

\renewcommand{\VestnikData}{01.11.09} % Дата подписания Вестника в печать

\renewcommand{\VestnikUPL}{26,64} % Усл. печ. л.

\renewcommand{\VestnikUIL}{25,73} % Уч.-изд. л.

\renewcommand{\VestnikRegNumber}{479/09} % Регистрационный номер

\renewcommand{\VestnikOrder}{994} % Номер заказа







%

%%%%%

%%%   Самарский государственный технический университет

%%%   Кафедра прикладной математики и информатики

%%%

%%%   Преамбула для статьи в журнал "Вестник Самарского государственного

%%%   технического университета. Серия: «Физико-математические науки»"

%%%   Ориентирована под MikTEx 2.4 for Win32

%%%

%%%   Хотя сам журнал собирается под GNU/Linux на дистрибутиве TeXLive



\documentclass[11pt,a4paper,twoside]{article}



\usepackage[cp1251]{inputenc}

\usepackage[T2A]{fontenc}

\usepackage[english,russian]{babel}

\usepackage{samgtubib}

\usepackage[dvips]{graphicx}

\usepackage[multiple]{footmisc}



\usepackage{samgtu}

\renewcommand{\VestnikYear}{2009} % Год издания Вестника

\renewcommand{\VestnikNumber}{2\,(19)} % Номер Вестника

\renewcommand{\VestnikMonth}{Ноябрь} % Месяц выхода Вестника

\renewcommand{\VestnikData}{01.11.09} % Дата подписания Вестника в печать

\renewcommand{\VestnikUPL}{26,64} % Усл. печ. л.

\renewcommand{\VestnikUIL}{25,73} % Уч.-изд. л.

\renewcommand{\VestnikRegNumber}{479/09} % Регистрационный номер

\renewcommand{\VestnikOrder}{994} % Номер заказа







%

%\usepackage{qrcode}

%

%%

%%\newcommand{\totop}[1]{\raisebox{6pt}[1pt][2pt]{#1}} 

%%\newcommand{\tottop}[1]{\raisebox{12pt}[1pt][2pt]{#1}} 

%%\newcommand{\totttop}[1]{\raisebox{18pt}[1pt][2pt]{#1}} 

%%\newcommand{\tottttop}[1]{\raisebox{24pt}[1pt][2pt]{#1}} 

%%\newcommand{\totttttop}[1]{\raisebox{30pt}[1pt][2pt]{#1}} 

%\newcommand{\tobot}[1]{\raisebox{-6pt}[1pt][2pt]{#1}} 

%%\newcommand{\tobbot}[1]{\raisebox{-12pt}[1pt][2pt]{#1}} 

%%\newcommand{\tobbbot}[1]{\raisebox{-18pt}[1pt][2pt]{#1}}

%

%\graphicspath{{./00PIC/}}

%

%%\samgtupubl % для генерации pdf, отдаваемого в типографию СамГТУ

%\pdfcrop % обрезка pdf для размещения на math-net.ru

%\draft % На корректуре

%\filllastpage % Последняя страница не заполнена не полностью

%%\briefcomm % Статья является кратким сообщением



\vsgtu{1485}



\FirstPage{220}

\LastPage{240}

%%

%%

%\begin{document}







\renewcommand{\baselinestretch}{0.9}



%\hfill \qrcode[hyperlink,height=.8in]{http://dx.doi.org/10.14498/vsgtu1485}

%\vspace{-.8in}



%%%%%%%%%%%%%%%%%%%%%%%%%%%%%%%%%%%%%%%%%%%%%%%%%%%%%%%%%%%%%%%%%%%%%%%%%%%%%%%%%%%%

%Информация о статье и об авторах на русском и английском языках





\udc{517.956.6}

\msc{35M12}% Коды MSC см. http://www.ams.org/msc/

\rutitle{Нелокальная задача  для нагруженного уравнения~смешанного типа с~интегральным~оператором}% Полный заголовок

{Нелокальная задача  для нагруженного уравнения смешанного типа\dots}% Заголовок, который идёт в верхний колонтитул



\ruauthor{О.~Х.~Абдуллаев}% Автор(ы) для заголовка, если авторов несколько укать \inst

{А\,б\,д\,у\,л\,л\,а\,е\,в~О.~Х.}% Автор(ы) для оглавления и колонтитула через \,





\ruaffil{\rusnuu.}

\enaffil{\engnuu.}



%\email{E-mails}{obidjon.mth@gmail.com} %E-mail's авторов obid_mth@yahoo.com



\ruauthordetails{1}{\fullauthor{\rupid{66335}{Обиджон Хайруллаевич Абдуллаев}} \regalia{к.ф.-м.н., доц.; \mem{obidjon.mth@gmail.com}} \position{доцент} \dept{каф. дифференциальных уравнений и математической физики}}





\enauthordetails{1}{\fullauthor{Obidjon Kh. Abdullayev} \regalia{Cand. Phys. \& Math. Sci.; \mem{obidjon.mth@gmail.com}}

\position{Associate Professor} \dept{Dept. of Differential Equations and Mathematical Physics}}



\rukeywords{нагруженное уравнение, интегральный оператор, уравнения эллиптико"=гиперболического типа, двусвязная область, существование и единственность решения, принцип экстремума, интегральные уравнения}



\enkeywords{loaded equation, integral operator, elliptic-hyperbolic type equations, double-connected domain, existence and uniqueness of solution, extremum principle, integral equations}



\ruabstract{Поставлена и исследована нелокальная задача для нагруженного уравнения второго порядка

эллиптико"=гиперболического типа  с~интегральным оператором в~двусвязной области.

Единственность решения доказывается с~помощью принципа экстремума для уравнений смешанного типа.

Для использования принципа экстремума было показано, что нагруженная часть уравнения тождественно равна нулю. 

Существование решения задачи доказывается методом интегральных уравнений, при этом используются теория сингулярных интегральных уравнений и~интегральные уравнения Фредгольма второго рода.}



%

\entitle{A non-local problem for a  loaded mixed-type equation  with a integral operator}



\enauthor{O.~Kh.~Abdullayev}{A\,b\,d\,u\,l\,l\,a\,y\,e\,v~O.~Kh.}



\enabstract{We study the existence and uniqueness of the solution

of non-local boundary value problem for the loaded elliptic-hyperbolic

equation

$$

u_{xx} + \mathop{\mathrm{sgn}} (y) u_{yy} + \frac{1 - \mathop{\mathrm{sgn}} (y)}{2}

\sum\limits_{k = 1}^n {R_k}(x, u(x, 0))  = 0 

$$

with integral operator

$$

{R_k}\bigl(x, u(x, 0)\bigr) = 

\left\{ 

\begin{array}{lc}

{p_k}(x)D_{x\,\,1}^{ - {\alpha _k}}u(x, 0), & q \le x \le 1,\\[2mm]

{r_k}(x)D_{ - 1\,x}^{ - {\beta _k}}u(x, 0), & - 1 \le x \le  - q, 

\end{array} \right.

$$

where 

$$

\begin{array}{l}

\displaystyle

D_{ax}^{ - {\alpha _k}}f(x) = \frac{1}{{\Gamma ({\alpha _k})}} \int _a^x \frac{f(t)}{(x - t)^{1-{\alpha _k} }}dt, 

\\

\displaystyle

 D_{xb}^{ - {\beta _k}}f(x) =

\frac{1}{{\Gamma ({\beta _k})}} \int _x^b \frac{f(t)}{(t - x)^{1-{\beta _k}}}dt ,

\end{array}

$$

 in double-connected domain   $\Omega $, bounded with two lines:

$$

 \sigma _1:~x^2 + y^2 = 1,\quad \sigma _2:~ x^2 + y^2 = q^2 \quad \text{at $y > 0$,}$$ 

 and characteristics:

$$ A_j C_1:~ x + ( - 1)^j y = ( - 1)^{j + 1},\quad B_j C_2:~x + ( - 1)^j y = ( - 1)^{j + 1} \cdot q$$

of the considered equation at $y < 0$, 

where $0 < q < 1$, $j = 1, 2$; $A_1 ( 1; 0),$ $A_2(  - 1; 0)$, 

$B_1(q; 0)$,  $B_2( - q; 0)$,  $C_1(0; - 1)$, $C_2(0; - q)$, $\beta _k$, $\alpha _k > 0$.



Uniqueness of the solution of investigated problem was proved by

an extremum principle for the mixed type equations. Thus we need

to prove that, the loaded part of the equation is identically

equal to zero if considerate problem is homogeneous. Existence

of the solution of the problem was proved by a method of the

integral equations, thus the theory of the singular integral

equations and Fredholm integral equations of the second kind were

widely used.}









\date{10/III/2016}

\revisiondate{25/IV/2016}

\accepteddate{27/V/2016}









\makerutitle



Первые фундаментальные исследования в теории нагруженных уравнений принадлежат А.~М.~Нахушеву \cite{Naxush:bib:Naxushev, Nax:bib:Naxushe}. 

В этих работах даны наиболее общее определение нагруженного уравнения и подробная классификация нагруженных уравнений: нагруженных дифференциальных, интегральных, интегро"=дифференциальных, функциональных  уравнений, а также их многочисленные приложения. 

За этими исследованиями последовали работы А.~М.~Нахушева \cite{Naxusha:bib:Naxushevam},   В.~А.~Елеева и~А.~В.~Дзарахохова~\cite{Eleev:bib:Eleeva, Elev:bib:Eleevdzrox}, В.~М.~Казиева~\cite{Kaz:bib:Kaziyev, Kaze:bib:Kaziev}, И.~Н.~Ланина \cite{Lan:bib:Lanin}, Б.~И.~Исломова и~Д.~М.~Курьязова~\cite{Islom:bib:Islomov, Kur:bib:Kuryaz}, М.~И.~Рамазанова \cite{Ram:bib:Ramaz} и~К.~У.~Хубиева~\cite{Xub:bib:Xubiyev}, посвященные краевым задачам для нагруженных уравнений гиперболического, параболического, эллиптического и смешанного типов второго порядка, в которых были получены интересные результаты.

В работах К.~Б.~Сабитова и Е.~П.~Мелишевой~[\citen{Sab:bib:Sabit}--\citen{Meli:bib:Melish}] исследованы разные задачи для нагруженного уравнения смешанного типа в~четырехугольных областях. 

Приведенные  публикации  позволяют сделать вывод, что вопросы теории нагруженных уравнений являются актуальными, а теория нагруженных уравнений находит свое место в приложениях и быстро развивается.

Однако краевые задачи для нагруженных уравнений смешанного типа  с интегральным оператором дробного порядка в~двусвязных областях до сих пор недостаточно исследованы. 

Отметим, что локальные и нелокальные краевые задачи для нагруженного уравнения эллиптико"=гиперболического типа в~двусвязной области были исследованы в~работах \cite{Abdu:bib:Abdull, Abd:bib:Abdullaev}, в~которых рассматриваются уравнения с~двумя линиями  изменения типа, при этом в~нагруженной части участвует не дифференциальный оператор, а просто след самой функции.

В настоящей работе рассматривается уравнение с~одной линией изменения типа, у которого в~нагруженной части присутствует дифференциальный оператор и используются нелокальные условия более общего вида.





\smallskip



%%%%%%%%%%%%%%%%%%%%%%%%%%%%%%%%%%%%%%%%%%%%%%%%%%%%%%%%%%%%%%%%%%%%%%%%%%%%%%%%%%%%

\Section{Постановка задачи} % Нумерованный раздел

В  конечной двусвязной области $\Omega $, описание которой приводится ниже, рассматривается 

%нелокальная задача для 

нагруженное уравнение эллиптико"=гиперболического типа

\begin{equation}

u_{xx} + \mathop{\rm sgn} (y)u_{yy} + \frac{{1 - \mathop{\rm sgn} (y)}}{2}\sum\limits_{k = 1}^n {R_k}\bigl(x, u(x, 0)\bigr)  = 0 \label{eq:1.1}

\end{equation}

c интегральным оператором

$$

{R_k}\bigl(x, u(x, 0)\bigr) = 

\left\{ 

\begin{array}{lc}

{p_k}(x)D_{x\,\,1}^{ - {\alpha _k}}u(x, 0), & q \le x \le 1,\\[2mm]

{r_k}(x)D_{ - 1\,x}^{ - {\beta _k}}u(x, 0), & - 1 \le x \le  - q, 

\end{array} \right.$$

где

\begin{equation}

\begin{array}{l}

\displaystyle

D_{ax}^{ - {\alpha _k}}f(x) = \frac{1}{{\Gamma ({\alpha _k})}} \int _a^x \frac{f(t)}{(x - t)^{1-{\alpha _k} }}dt, 

\\[2mm]

\displaystyle

 D_{xb}^{ - {\beta _k}}f(x) =

\frac{1}{{\Gamma ({\beta _k})}} \int _x^b \frac{f(t)}{(t - x)^{1-{\beta _k}}}dt ,

\end{array}

\label{eq:1.2}

\end{equation}

${\beta _k},$ ${\alpha _k}$ "--- положительные константы.

Область $\Omega $ при $y > 0$ ограничена следующими линиями:

 $${\sigma _1}:\,\,{x^2} + {y^2} = 1,\quad {\sigma _2}:\,\,{x^2} + {y^2} = {q^2},$$ 

 при $y < 0$ "--- следующими характеристиками уравнения \eqref{eq:1.1}:

$${A_j}{C_1}:\, x + {( - 1)^j}y = {( - 1)^{j + 1}}, \quad {B_j}{C_2}:\,x + {( - 1)^j}y = {( - 1)^{j + 1}} \cdot q ,

$$

где 

$0 < q < 1$, $j = 1, 2$;

${A_1}( 1; 0),$ $A_2 (  - 1; 0),$ $B_1( q; 0),$  $B_2(  - q; 0),$  $C_1( 0; - 1),$ $C_2( 0; - q)$.



 Введем следующие обозначения: 

 $${\theta _1}(x) = \frac{{x + 1}}{2} + i \cdot \frac{{x - 1}}{2},\quad 

 {\theta _2}(x) = \frac{{x - 1}}{2} - i \cdot \frac{{x +  1}}{2},\quad {i^2} = - 1;$$

$${\Omega _0} = \Omega  \cap (y > 0),\quad {\Delta _1} = \Omega  \cap (x + y > q) \cap (y < 0),$$

$${\Delta _2} = \Omega  \cap (y - x > q) \cap (y < 0),\quad {D_1} = \Omega  \cap ( - q < x + y < q) \cap (x > 0),$$

$${D_2} = \Omega  \cap ( - q < y - x < q)  \cap (x < 0),$$

$${D_3} = \Omega  \cap ( - 1 < x + y <  - q) \cap ( - 1 < y - x <  - q),$$

$${I_{2 + j}} = \bigl\{ {x:\,0 < {{( - 1)}^{j - 1}}x < q} \bigr\},\quad 

{I_j} = \Bigl\{ {x:\frac{{q + 1}}{2} < {{( - 1)}^{j - 1}}x < 1} \Bigr\}, \quad j = 1, 2.$$



В области $\Omega $ для уравнения \eqref{eq:1.1} ставится и исследуется следующая нелокальная задача.



\smallskip







\hypertarget{abd:taskI}{}



{\small \sc Задача I.} 

\it 

Найти функцию $u(x,y)$ со следующими свойствами{\rm\/:}

\begin{itemize}

\item[\rm 1)] $u(x,y) \in C(\bar \Omega );$

\item[\rm 2)] $u(x, y)$ является регулярным решением уравнения \eqref{eq:1.1} в области 

$$\Omega \backslash (y - x =  \pm q)\backslash (x + y =  \pm q),$$ кроме того$,$ ${u_y} \in C({A_1}{B_1} \cup A_2B_2),$ причем ${u_y}(x, - 0)$ и

${u_y}(x, + 0)$ могут обращаться в бесконечность порядка меньше единицы при $x \to  \pm q,$ а~при  $x \to  \pm 1$ ограничены$;$

\item[\rm 3)]  на линиях изменения типа выполняются условия сопряжения

\begin{equation}

{u_y}(x, + 0) = {\lambda _j}(x){u_y}(x, - 0),\quad (x, 0) \in {A_j}{B_j}, 

\label{eq:3.j}

\end{equation}

\item[\rm 4)]   $u(x,y)$ удовлетворяет краевым условиям

\begin{gather}

u(x,y)\bigr|_{\sigma _j}  = {\varphi _j}(x,y),\quad (x,y) \in \overline {\sigma _j};  

\label{eq:4.j}

\\

u(x,y)\bigr|_{B_jC_2} = g_{_j}(x), \quad x \in \overline {I_{2 + j}}; 

\quad  j=1, 2 \label{eq:5.j}

\end{gather}

\begin{multline}

\frac{d}{{dx}}u \bigl( \theta _1(x)\bigr) 

= {a_1}(x){u_y}(x, 0) + {b_1}(x){u_x}(x, 0) + \\ + {c_1}(x)u(x, 0) + {d_1}(x),\quad  x \in (q, 1),\label{eq:6.1}

\end{multline}

\begin{multline}

\frac{d}{{dx}}u\bigl(\theta _2(x)\bigr) 

= {a_2}(x){u_y}(x, 0) + {b_2}(x){u_x}(x, 0) + \\ + {c_2}(x)u(x, 0) + {d_2}(x),\quad  x \in (- 1,- q),\label{eq:6.2}

\end{multline}

где ${\varphi _j}(x,y),$ ${\lambda _j}(x),$ ${g_j}(x),$ ${a_j}(x),$ ${b_j}(x),$ ${c_j}(x),$ ${d_j}(x)$ 

$(j = 1, 2)$ и $p_k(x),$ $r_k (x)$ "--- заданные функции, причем

\begin{equation}

{g_1}( 0 ) = {g_2}(0),\quad 

{g_2}(  - q) = {\varphi _2}( - q, 0),\quad 

{g_1}( q ) = {\varphi _2}(q, 0).

\label{eq:1.7}

\end{equation}





\end{itemize}



\rm 



\smallskip



{\sc\small Лемма.} 

\it 

Общее решение уравнения \eqref{eq:1.1} при $y < 0$ представимо в виде

\begin{equation}

u(x,y) = {f_1}(x + y) - {f_2}(x - y) + \omega (x),

\label{eq:1.8}

\end{equation}

где ${f_1}(x,y)$ и ${f_2}(x,y)$ "--- произвольные дважды непрерывно дифференцируемые функции$,$ 

$$\omega (x) = 

\left\{ \begin{array}{ll}

{\omega _1}(x), & x \ge 0,\\

{\omega _2}(x), & x \le 0,

\end{array} \right.$$  

причем

\begin{gather}

{\omega _1}(x) =  -  \int _x^1 {dt \int _t^1 {\sum\limits_{k = 1}^n {{p_k}(z)D_{z\,\,1}^{ - {\alpha _k}}u(z, 0)} d} } z ; \label{eq:9.1}

\\

{\omega _2}(x) =  -  \int _{ - 1}^x {dt \int _{ - 1}^z {\sum\limits_{k = 1}^n {{r_k}(z)D_{ - 1\,z}^{ - {\beta _k}}u(z, 0)} d} } z. \label{eq:9.2}

\end{gather}



\rm



\smallskip



\begin{proof}

Пусть функция \eqref{eq:1.8} является решением уравнения \eqref{eq:1.1}, 

тогда, подставив \eqref{eq:1.8} в \eqref{eq:1.1}, получим

 \begin{equation}

 \omega ''(x) + \sum\limits_{k = 1}^n {R_k} \bigl(x, u(x, 0) \bigr)  = 0.

 \label{eq:1.10}

 \end{equation}

Очевидно, что функции \eqref{eq:9.1} и \eqref{eq:9.2} соответственно при $x \ge 0$ и $x \le 0$ удовлетворяют уравнению \eqref{eq:1.10}.

Теперь, наоборот, пусть $u(x,y)$  "--- регулярное решение уравнения \eqref{eq:1.1}. 

Докажем, что функция \eqref{eq:1.8} является общим решением уравнения~\eqref{eq:1.1}.



Нетрудно проверить, что функции  \eqref{eq:9.1} и \eqref{eq:9.2} являются частными решениями уравнения 

$${u_{xx}} - {u_{yy}} + \sum\limits_{k = 1}^n {R_k} \bigl(x,u(x, 0)\bigr)  = 0$$ соответственно при $x \ge 0$ и $x \le 0$, а

функция ${f_1}(x + y) - {f_2}(x - y)$ является общим решением уравнения ${u_{xx}} - {u_{yy}} = 0$.



Следовательно, функция \eqref{eq:1.8} является общим решением уравнения \eqref{eq:1.1}. 

\end{proof}



\smallskip



\Section {Единственность решения задачи \hyperlink{abd:taskI}{I}}

В~силу  \eqref{eq:1.8} и \eqref{eq:9.1}, \eqref{eq:9.2} утверждаем, что решение задачи Коши для уравнения \eqref{eq:1.1} в~области ${\Delta _j}$ ($j = 1, 2$),

удовлетворяющее условиям   

$$

u(x, 0) = {\tau_j}(x),\quad  x\in \overline {{A_j}{B_j}} $$ и  

$${u_y}(x, - 0) = \nu _j^ - (x),\quad x \in {A_j}{B_j},$$  

имеет следующий вид:

 \begin{multline}

u(x,y) = \frac12 \bigl({{{\tau _1}(x + y) + {\tau _1}(x - y)}} \bigr) -

 \frac{1}{2} \int _{x + y}^{x - y} {\nu _1^ - (t)dt}  +

 \\ 

 +

  \frac{1}{2} \int _{x + y}^1

{dt \int _t^1 {\sum\limits_{k = 1}^n {{p_k}(z)D_{z\,\,1}^{ - {\alpha _k}}{\tau _1}(z)} dz} } +

 \frac{1}{2} \int _{x - y}^1 {dt \int _t^1 {\sum\limits_{k = 1}^n {{p_k}(z)D_{z\,\,1}^{- {\alpha _k}}{\tau _1}(z)}dz}} - 

 \\-  \int _x^1

 {dt \int _t^1 {\sum\limits_{k = 1}^n {{p_k}(z)D_{z\,\,1}^{ - {\alpha _k}}u(z, 0)} d} }z, \label{eq:11.1}

\end{multline}

\begin{multline}

u(x,y) = \frac12 \bigl({{{\tau _2}(x + y) + {\tau _2}(x - y)}}\bigr)  - \frac{1}{2} \int _{x + y}^{x - y} {\nu _2^ - (t)dt}  +

\\

+ \frac{1}{2} \int _{ - 1}^{x + y} {dt \int _{ - 1}^t {\sum\limits_{k = 1}^n {{r_k}(z)D_{ - 1\,z}^{ - {\beta _k}}{\tau _2}(z)} dz} }

+ \frac{1}{2} \int _{ - 1}^{x - y} {dt \int _{ - 1}^t {\sum\limits_{k = 1}^n {{r_k}(z)D_{ - 1\,z}^{ - {\beta _k}}{\tau _2}(z)} dz} }  - \\ -

  \int _{ - 1}^x {dt \int _{ - 1}^z {\sum\limits_{k = 1}^n {{r_k}(z)D_{ - 1\,z}^{ - {\beta _k}}u(z, 0)} d} } z.\label{eq:11.2}

\end{multline}

Отсюда, учитывая, что

\begin{multline}

u\bigl( {{\theta _1}(x)} \bigr) = \frac12 \bigl({{{\tau _1}(x) + {\tau _1}(1)}}\bigr)  - \frac{1}{2} \int _x^1 {\nu _1^ - (t)dt}  + \\

+ \frac{1}{2} \int _x^1

{dt \int _t^1 {\sum\limits_{k = 1}^n {{p_k}(z)} } D_{z\,1}^{ - {\alpha _k}}{\tau _1}(z)dz}  -

\\

-  \int _{({{x + 1}})/{2}}^1

{dt \int _t^1 {\sum\limits_{k = 1}^n {{p_k}(z)D_{z\,\,1}^{ - {\alpha _k}}u(z, 0)} d} } z,

\end{multline}

\begin{multline}

u\bigl( {{\theta _2}(x)} \bigr) = \frac12 \bigl({{{\tau _2}( - 1) + {\tau _2}(x)}}\bigr) - \frac{1}{2} \int _{ - 1}^x {\nu _2^ - (t)dt}  +

\\

+

\frac{1}{2} \int _{ - 1}^x {dt \int _{ - 1}^t {\sum\limits_{k = 1}^n {{r_k}(z)} } D_{ - 1\,z}^{ - {\beta _k}}{\tau _2}(z)dz} -

\\ 

-  \int _{

- 1}^{({{x - 1}})/{2}} {dt \int _{ - 1}^t {\sum\limits_{k = 1}^n {{r_k}(z)D_{ - 1\,z}^{ - {\beta _k}}u(z, 0)} d} } z,

\end{multline}

в силу условий \eqref{eq:6.1} и \eqref{eq:6.2}  получим

\begin{multline}

\frac{1}{2}{\tau '_1}(x) + \frac{1}{2}\nu _1^ - (x) - \frac{1}{2}\sum\limits_{k = 1}^n { \int _x^1 {{p_k}(z)D_{z1}^{ - {\alpha _k}}{\tau _1}(z)dz} } +

\\

+ \frac{1}{2}\sum\limits_{k = 1}^n { \int _{({{x + 1}})/{2}}^1 {{p_k}(z)D_{z1}^{ - {\alpha _k}}{\tau _1}(z)dz} }  = 

\\

 = {a_1}(x)\nu _1^ - (x) + {b_1}(x){\tau '_1}(x) + {c_1}(x){\tau _1}(x) + {d_1}(x), \label{eq:12.1}

\end{multline}

\begin{multline}

\frac{1}{2}{\tau '_2}(x) - \frac{1}{2}\nu _2^ - (x) + \frac{1}{2}\sum\limits_{k = 1}^n { \int _{ - 1}^x {{r_k}(z)D_{ - 1\,z}^{ - {\beta _k}}{\tau _2}(z)dz} } - \\

  - \frac{1}{2}\sum\limits_{k = 1}^n { \int _{ - 1}^{({{x - 1}})/{2}} {{r_k}(z)D_{ - 1\,z}^{ - {\beta _k}}{\tau _2}(z)dz} }  =

\\

= {a_2}(x)\nu _2^ - (x) + {b_2}(x){\tau '_2}(x) + {c_2}(x){\tau _2}(x) + {d_2}(x).

\label{eq:12.2}

\end{multline}

\pagebreak



\noindent

Следовательно,

\begin{multline}

 \bigl( {2{a_1}(x) - 1} \bigr)\nu _1^ - (x) = 

\\

= 

 \bigl( {1 - 2{b_1}(x)} \bigr){\tau '_1}(x) - 2{c_1}(x){\tau _1}(x) 

 - \frac{1}{2}\sum\limits_{k = 1}^n { \int _x^1 {{p_k}(z)D_{z\,1}^{ - {\alpha _k}}{\tau _1}(z)dz} } +

 \\

  + \frac{1}{2}\sum\limits_{k = 1}^n

 { \int _{({{x + 1}})/{2}}^1 {{p_k}(z)D_{z\,1}^{ - {\alpha _k}}{\tau _1}(z)dz} }  - 2{d_1}(x),

 \label{eq:13.1}

\end{multline}

\begin{multline}

\bigl(2{a_2}(x) + 1\bigr)\nu _2^ - (x) = \\

= \bigl(1 - 2{b_2}(x)\bigr){\tau '_2}(x) - 2{c_2}(x){\tau _2}(x) 

+ \frac{1}{2}\sum\limits_{k = 1}^n { \int _{ - 1}^x {{r_k}(z)D_{ - 1\,z}^{ - {\beta _k}}{\tau _2}(z)dz} }  -

\\

 - \frac{1}{2}\sum\limits_{k = 1}^n

{ \int _{ - 1}^{({{x - 1}})/{2}} {{r_k}(z)D_{ - 1\,z}^{ - {\beta _k}}{\tau _2}(z)dz} }  - 2{d_2}(x).\label{eq:13.2}

\end{multline}



\smallskip



\hypertarget{abd:th1}{}



{\small\sc Теорема 1.} \it  Если ${\beta _k},$ ${\alpha _k} > 0$ и выполнены условия

\begin{gather}

2{a_j}(x) + {( - 1)^j} > 0,\quad  {c_j}(x) \le 0,\quad  {\lambda _j}(x) > 0,\quad j = 1, 2;

\label{eq:14.j}

\\

{p_k}(x) > 0,\quad {r_k}(x) > 0,\quad k = 1, 2, \dots, n, 

\label{eq:1.15}

\end{gather}

то решение $u(x,y)$ задачи~\hyperlink{abd:taskI}{\rm I} единственно\rm.



\rm



\smallskip



\begin{proof}

Известно, что если решение однородной задачи  (т.~е. решение однородного уравнения с~однородными условиями) является тождественным нулем, то решение соответствующей неоднородной задачи единственно.   

Следовательно, достаточно показать, что решение $u(x,y)$ задачи~\hyperlink{abd:taskI}{\rm I}   является тождественным нулем при ${\varphi _j}(x, y) \equiv {d_j}(x) \equiv 0$.



В~силу принципа экстремума для эллиптических уравнений решение $u(x,y)$ уравнения \eqref{eq:1.1} достигает своих экстремальных значений лишь на границе области ${\bar \Omega _0}$, 

т.~е. на ${\bar \sigma _1} \cup {\bar \sigma _1} \cup \overline {{A_1}{B_1}}  \cup \overline {{A_2}{B_2}} $.



 Пусть ${\varphi _j}(x,y) \equiv 0$, тогда учитывая, что $u(x,y) \in C(\bar \Omega )$,  имеем 

 $${\tau _1}(1) = {\varphi _1}(1, 0) =  0 \quad \text{и}\quad {\tau _1}(q) = {\varphi _2}(q, 0) = 0.$$ 

Далее предположим, что функция ${\tau _1}(x)$  обращается в нуль хотя бы в одной точке интервала $(q, 1)$, т.~е. предположим, что  ${x_i}$ ($i = \overline {1, m} $) "--- нули функции~${\tau _1}(x)$.



Рассмотрим отрезок $[ {x_1}, {x_2}] \subset \overline {{A_1}{B_1}}$. 

Так как ${\tau _1}({x_1}) = {\tau _1}( {x_2}) = 0,$ 

функция ${\tau _1}( x ) > 0$ или ${\tau _1}( x) < 0$ для всех $x \in ( x_1 , x_2)$. 



Предположим ${\tau _1}\left( x \right) > 0$ $\left( {{\tau _1}\left( x \right) < 0} \right),$ тогда можно показать, 

что внутри  этого интервала функция ${\tau _1}(x)$ не достигает своего положительного максимума (отрицательного минимума).

Действительно, если в точке ${x_0} \in ( x_1, x_2)$ функция ${\tau _1}(x)$ достигает своего положительного максимума (отрицательного минимума), то из \eqref{eq:13.1} при ${d_1}(x) \equiv 0$ получим

\begin{multline}

\bigl( {2{a_1}({x_0}) - 1} \bigr)\nu _1^ - ({x_0}) =  - 2{c_1}({x_0}){\tau _1}({x_0}) + 

\\

+ \frac{1}{2}\sum\limits_{k = 1}^n 

\biggl( \int _{({{{x_0} + 1}})/{2}}^1 {{p_k}(z)D_{z\,1}^{ - {\alpha _k}}{\tau _1}(z)dz}  -

  \int _{{x_0}}^1 {{p_k}(z)D_{z\,1}^{ - {\alpha _k}}{\tau _1}(z)dz}  

  \biggr). 

\end{multline}



Отсюда, учитывая,  что $D_{x\,1}^{ - {\alpha _k}}{\tau _1}(x) > 0$ ($D_{x\,1}^{ - {\alpha _k}}{\tau _1}(x) < 0$) при \mbox{${\tau _1}(x) > 0$} \mbox{(${\tau _1}(x) < 0$)}, в силу \eqref{eq:14.j} (при $j=1$), \eqref{eq:1.15}    получим, 

что $\nu _1^ - ({x_0}) \ge 0$ $\left( {\nu _1^ - ({x_0}) \le 0} \right)$, 

а~это  в~силу \eqref{eq:3.j} (учитывая, что ${\lambda _j}(x) > 0$) противоречит известному принципу Заремба"--~Жиро \cite{Bits:bib:Bitsad}, согласно которому в~точке положительного максимума (отрицательного минимума) должно выполняться условие

${\nu _1}({x_0}) < 0$ $\left( {{\nu _1}({x_0}) > 0} \right)$. 

Следовательно, ${\tau _1}(x)$

не  достигает своего положительного максимума (отрицательного минимума) в точке ${x_0} \in  ( {x_1}, {x_2})$.



Таким образом,

\begin{equation}

{\tau _1}(x) = 0 \quad  \forall x \in [ {x_1}, {x_2}].\label{eq:1.16}

\end{equation}

Аналогично вышеизложенным методом доказывается, что ${\tau _1}(x)$ не  достигает своего положительного максимума и отрицательного минимума { и в других интервалах},

т.~е.

\begin{equation}

\tau_1(x) = 0. \label{eq:1.17}

\end{equation}





Если ${x_1} = q$, т.~е. $$\tau ({x_1}) = {\tau _1}(q) = {\varphi _2}(q, 0) = 0,$$

то из \eqref{eq:1.16} и \eqref{eq:1.17} следует, что 

${\tau _1}(x) = 0$ $\forall x \in {A_1}{B_1},$ 

а~при ${x_1} \ne q$ функция ${\tau _1}(x)$ не может иметь экстремума в интервале $( q, {x_1})$, 

т.~е. функция ${\tau _1}(x)$ либо знакопостоянна в $(q, x_{1})$, 

либо  ${\tau _1}(x) = 0$ $\forall x \in [ q, x_{1} ]$.  



Предположим, что

функция ${\tau _1}(x)$ знакопостоянна в $(q,\,x_{1}),$ тогда, принимая во внимание, что ${\tau _1}({x_1}) = {\tau _1}(q) = 0$, заключаем, что

\begin{equation}

{\tau _1}(x) = 0\quad \forall x \in [ q, x_1]. \label{eq:18.1}

\end{equation}



Точно так же, если ${x_m} = 1$, т.~е. $${\tau _1}({x_m}) = {\tau _1}(1) = {\varphi _1}(1, 0) = 0,$$ 

то из \eqref{eq:1.16} и \eqref{eq:1.17} следует, что 

${\tau _1}(x) = 0$ $\forall x \in {A_1}{B_1}$, а при ${x_m} \ne 1$ функция ${\tau _1}(x)$ не может иметь экстремума в~интервале $({x_m}, 1)$, т.~е.

функция ${\tau _1}(x)$ либо знакопостоянна в ${x_m} \ne 1$ , либо ${\tau _1}(x) = 0$ $\forall x \in [ {x_m}, 1]$. 



Предположим, что функция ${\tau_1}(x)$ знакопостоянна в $({x_m}, 1),$ тогда, принимая во внимание, что 

$${\tau _1}({x_m}) = {\tau _1}(1) = 0,$$ 

заключаем, что 

\begin{equation}

{\tau _1}(x) = 0\quad \forall x \in [ x_m, 1].

\label{eq:18.2}

\end{equation}

В силу \eqref{eq:1.17} и \eqref{eq:18.1}, \eqref{eq:18.2} имеем ${\tau _1}(x) = 0$ $\forall x \in [q, 1]$.

Аналогичным  образом доказывается, что ${\tau _2}(x) = 0$ $\forall x \in {A_2}{B_2}$.



Таким образом, показано, что  ${\tau _j}(x) \equiv 0$ 

$\forall x \in \overline {{A_j}{B_j}} $ при ${\varphi _j}(x,y) \equiv {d_j}(x) \equiv 0$ $(j= 1, 2)$.

Кроме этого, показано, что решение $u(x, y)$ не достигает своего положительного максимума (отрицательного минимума) в интервалах

$(A_1, B_1)$ и $(A_2,B_2).$

Следовательно, учитывая, что $u(x,y) \in C\left( {{{\overline \Omega  }_0}} \right)$, на основании краевых условий  \eqref{eq:4.j}  при  ${\varphi _j}(x,y) \equiv 0$ $(j = 1, 2)$ имеем, что $u(x,y) \equiv 0$ в замкнутой области ${\overline \Omega  _0}$.



В силу того, что ${\tau _1}(x) \equiv {\tau _2}(x) \equiv 0$, из функциональных соотношений \eqref{eq:13.1} и \eqref{eq:13.2} получим, что ${\nu _1}(x) \equiv {\nu _2}(x) \equiv 0$. 

Следовательно, решение задачи Коши для уравнения \eqref{eq:1.1} в области ${\Delta _j}$ тождественно равно нулю, т.~е. $u(x,y)

\equiv 0$ в областях ${\Delta _j}$. 

В~силу единственности решения задачи Гурса имеем, что $u(x,y) \equiv 0$ в областях ${D_j}$ $(j = 1, 2, 3)$. 

Таким образом, $u(x,y) \equiv 0$ в области $\overline \Omega  $ и тем самым доказана единственность решения задачи~\hyperlink{abd:taskI}{\rm I}.

\end{proof}



\smallskip



\Section{Существование  решения задачи~\hyperlink{abd:taskI}{\rm I}}



\hypertarget{abd:th2}{}



{\small\sc Теорема~2.} \it 

Если выполнены условия теоремы~\hyperlink{abd:th1}{\rm 1} и условия

\begin{equation}

{\varphi _j}(x,y) = {(xy)^\gamma }\overline {{\varphi _j}} (x,y), \quad \overline {{\varphi _j}} (x,y) \in C\left( {\overline {{\sigma _j}} } \right),\,\ 2

< \gamma  < 3, \label{eq:19.j}

\end{equation}

\begin{equation}

{a_j}(x),\; {b_j}(x),\; {c_j}(x),\; {d_j}(x),\; {\lambda _j}(x) \in C\left( {\overline {{I_j}} } \right) \cap {C^2}\left( {{I_j}} \right),  \label{eq:20.j}

\end{equation}

то решение задачи~\hyperlink{abd:taskI}{\rm I} существует.



\smallskip \rm



Для доказательства теоремы~\hyperlink{abd:th2}{\rm 2} нам потребуются следующие соотношения, которые получаются из \eqref{eq:13.1} и \eqref{eq:13.2} в силу \eqref{eq:1.2}:

\begin{multline}

\bigl( {1 - 2{b_1}(x)} \bigr){\tau '_1}(x) - 2{c_1}(x){\tau _1}(x) - \\

- \frac{1}{2}\sum\limits_{k = 1}^n {\frac{1}{{\Gamma ({\alpha _k})}} \int _x^1 {{p_k}(z)dz \int _z^1 {{{(t - z)}^{{\alpha _k} - 1}}} {\tau _1}(t)} dt}  

= \bigl( {2{a_1}(x) - 1} \bigr)\nu _1^ - (x) -

\\ - \frac{1}{2}\sum\limits_{k = 1}^n {\frac{1}{{\Gamma ({\alpha _k})}} \int _{({{x + 1}})/{2}}^1

 {{p_k}(z)dz \int _z^1 {{{(t - z)}^{{\alpha _k} - 1}}} {\tau _1}(t)} dt}  + 2{d_1}(x),

 \end{multline}

 \begin{multline}

\bigl(1 - 2{b_2}(x)\bigr){\tau '_2}(x) - 2{c_2}(x){\tau _2}(x) + \\

+ \frac{1}{2}\sum\limits_{k = 1}^n {\frac{1}{{\Gamma ({\beta _k})}} \int _{ - 1}^x

{{r_k}(z)dz \int _{ - 1}^z {{{(z - t)}^{{\beta _k} - 1}}} {\tau _2}(t)dt} }  

 = \bigl(2{a_2}(x) + 1\bigr)\nu _2^ - (x) +\\

  + \frac{1}{2}\sum\limits_{k = 1}^n {\frac{1}{{\Gamma ({\beta _k})}} \int _{ - 1}^{({{x - 1}})/{2}}

 {{r_k}(z)dz \int _{ - 1}^z {{{(z - t)}^{{\beta _k} - 1}}} {\tau _2}(t)dt} }  + 2{d_2}(x),

 \end{multline} 

или

 \begin{multline}

\bigl( {1 - 2{b_1}(x)} \bigr){\tau '_1}(x) + \frac{1}{2}\sum\limits_{k = 1}^n {\frac{1}{{\Gamma (1 +

{\alpha _k})}}}  \int _x^1 {{p_k}(z)dz}  \int _z^1 {{{(t - z)}^{{\alpha _k}}}} {\tau '_1}(t)dt -

\\

 - \frac{1}{2}\sum\limits_{k = 1}^n {\frac{1}{{\Gamma (1 + {\alpha _k})}} \int _{({{x + 1}})/{2}}^1 {{p_k}(z)dz \int _z^1 {{{(t - z)}^{{\alpha

 _k}}}} {\tau '_1}(t)} dt} 

  + 2{c_1}(x) \int _x^1 {{\tau '_1}(t)dt}= \\ =

\frac{1}{2}\sum\limits_{k = 1}^n {\frac{{{\tau _1}(1)}}{{\Gamma (1 + {\alpha _k})}} \int _x^1 {{p_k}(z){{(1 -

 z)}^{{\alpha _k}}}dz} }- \hspace{3cm}

\\- 

  \frac{1}{2}\sum\limits_{k = 1}^n {\frac{{{\tau _1}(1)}}{{\Gamma (1 + {\alpha _k})}} \int _{({{x + 1}})/{2}}^1 {{{(1 - z)}^{{\alpha

 _k}}}{p_k}(z)dz} } + 

 \\

  + 2{c_1}(x){\tau _1}(1) + 2{d_1}(x) + \bigl( {2{a_1}(x) - 1} \bigr)\nu _1^ - (x),\label{eq:21.1}

\end{multline}

 \begin{multline}

\bigl(1 - 2{b_2}(x)\bigr){\tau '_2}(x) + \frac{1}{2}\sum\limits_{k = 1}^n {\frac{1}{{\Gamma (1 + {\beta

_k})}} \int _{ - 1}^x {{r_k}(z)dz \int _{ - 1}^z {{{(z - t)}^{{\beta _k}}}} {\tau '_2}(t)dt} }-

\\

- \frac{1}{2}\sum\limits_{k = 1}^n {\frac{1}{{\Gamma (1 + {\beta _k})}} \int _{ - 1}^{({{x - 1}})/{2}} {{r_k}(z)dz \int _{ - 1}^z {{{(z -

 t)}^{{\beta _k}}}} {\tau '_2}(t)dt}}- 2{c_2}(x) \int _{ - 1}^x {{\tau '_2}(t)} dt  =

\\ 

 = \frac{1}{2}\sum\limits_{k = 1}^n {\frac{{{\tau _2}( - 1)}}{{\Gamma (1 + {\beta _k})}} \int _{ - 1}^{({{x - 1}})/{2}} {{{(z + 1)}^{{\beta _k}}}{r_k}(z)dz} } - \hspace{3cm}

\\ 

  - \frac{1}{2}\sum\limits_{k = 1}^n {\frac{{{\tau _2}( - 1)}}{{\Gamma (1 + {\beta _k})}} \int _{ - 1}^x {{r_k}(z){{(z + 1)}^{{\beta _k}}}dz} } -

  \\

 -2{c_2}(x){\tau _2}( - 1) + \bigl(2{a_2}(x) + 1\bigr)\nu _2^ - (x) + 2{d_2}(x).\label{eq:21.2}

\end{multline}



\smallskip



\Section{Основные функциональные соотношения}

При исследовании существования решения поставленной задачи рассматриваются следующие случаи.



\smallskip



\hypertarget{abd:A}{}



\texttt{Случай A.}

Пусть ${b_j}(x) \ne  - {1}/{2},$ тогда из \eqref{eq:21.1} и \eqref{eq:21.2} получим интегральные уравнения Вольтерры второго рода относительно ${\tau '_j}(x)$:

\begin{multline}

{\tau '_1}(x) +  \int _x^1 {\tau '_1}(t){K_{11}}(x,t)dt +  

\\

+ \int _{({{x + 1}})/{2}}^1 {{\tau '_1}(t){K_{12}}(x,t)dt}  = f_1(x), \quad  q \le x \le 1, \label{eq:22.1}

\end{multline}

\begin{multline}

{\tau '_2}(x) +  \int _{ - 1}^x {{\tau '_2}(t){K_{21}}(x,t)dt}  + \\ 

+ \int _{ - 1}^{({{x - 1}})/{2}} {{\tau '_2}(t){K_{22}}(x,t)dt}  = f_2(x),\quad      - 1 \le x \le  - q, \label{eq:22.2}

\end{multline}

где

\begin{equation}

{K_{11}}(x,t) = \frac{{2{c_1}(x)}}{{1 - 2{b_1}(x)}} + \frac{1}{{2\left( {1 - 2{b_1}(x)} \right)}}\sum\limits_{k = 1}^n {\frac{1}{{\Gamma (1 + {\alpha _k})}} \int _x^t {{{(t - z)}^{{\alpha _k}}}{p_k}(z)dz} },\label{eq:23.1}

\end{equation}

\begin{equation}

{K_{12}}(x,t) =  - \frac{1}{{2\left( {1 - 2{b_1}(x)} \right)}}\sum\limits_{k = 1}^n {\frac{1}{{\Gamma (1 + {\alpha _k})}} \int _{({{x + 1}})/{2}}^t {{{(t - z)}^{{\alpha _k}}}{p_k}(z)dz} }, \label{eq:23.2}

\end{equation}

\begin{equation}

{K_{21}}(x,t) = \frac{{2{c_2}(x)}}{{2{b_2}(x) - 1}} + \frac{1}{{2\left( {1 - 2{b_2}(x)} \right)}}\sum\limits_{k = 1}^n {\frac{1}{{\Gamma (1 + {\beta _k})}} \int _t^x {{{(z - t)}^{{\beta _k}}}{r_k}(z)dz} }, \label{eq:24.1}

\end{equation}

\begin{equation}

{K_{22}}(x,t) =  - \frac{1}{{2\left( {1 - 2{b_2}(x)} \right)}}\sum\limits_{k = 1}^n {\frac{1}{{\Gamma (1 + {\beta _k})}} \int _t^{({{x - 1}})/{2}} {{{(z - t)}^{{\beta _k}}}{r_k}(z)dz} }, \label{eq:24.2}

\end{equation}

\begin{multline}

{f_1}(x) = \frac{{2{a_1}(x) - 1}}{{1 - 2{b_1}(x)}}\nu _1^ - (x) - \frac{1}{2}\sum\limits_{k = 1}^n {\frac{{{\varphi _1}(1, 0)}}{{\Gamma (1 + {\alpha _k})}} \int _{({{x + 1}})/{2}}^1 {\frac{{{p_k}(z){{(1 - z)}^{{\alpha _k}}}}}{{1 - 2{b_1}(x)}}dz} }  + 

\\

+ \frac{1}{2}\sum\limits_{k = 1}^n {\frac{{{\varphi _1}(1, 0)}}{{\Gamma (1 + {\alpha _k})}} \int _x^1 {\frac{{{p_k}(z){{(1 - z)}^{{\alpha _k}}}}}{{1 - 2{b_1}(x)}}dz} }  + \frac{{2({d_1}(x) + {c_1}(x){\varphi _1}(1, 0))}}{{1 - 2{b_1}(x)}}, \label{eq:25.1}

\end{multline}

\begin{multline}

{f_2}(x) = \frac{{2{a_2}(x) + 1}}{{1 - 2{b_2}(x)}}\nu _2^ - (x) + \frac{1}{2}\sum\limits_{k = 1}^n {\frac{{{\varphi _1}( - 1, 0)}}{{\Gamma (1 + {\beta _k})}} \int _{ - 1}^{({{x - 1}})/{2}} {\frac{{{r_k}(z){{(1 + z)}^{{\beta _k}}}}}{{1 - 2{b_2}(x)}}dz} }  - \\

- \frac{1}{2}\sum\limits_{k = 1}^n {\frac{{{\varphi _1}( - 1, 0)}}{{\Gamma (1 + {\beta _k})}} \int _{ - 1}^x {\frac{{{r_k}(z){{(1 + z)}^{{\beta _k}}}}}{{1 - 2{b_2}(x)}}dz} }  + \frac{{2({d_2}(x) - {c_2}(x){\varphi _1}( - 1, 0))}}{{1 - 2{b_2}(x)}}. \label{eq:25.2}

\end{multline}

Учитывая, что ${b_j}(x) \ne {1}/{2}$, в силу \eqref{eq:19.j}, \eqref{eq:20.j}  из \eqref{eq:23.1}, \eqref{eq:23.2}, \eqref{eq:24.1} и \eqref{eq:24.2} получим 

$$\left| {{K_{1j}}(x,t)} \right| \le {c_1},\quad \left| {{K_{2j}}(x,t)} \right| \le {c_2},\quad  j = 1, 2.$$  

Учитывая класс функций $\nu _j^ - (x)$, в силу \eqref{eq:19.j}, \eqref{eq:20.j} из \eqref{eq:25.1} и \eqref{eq:25.2} получим (см. \cite{Islabd:bib:Islomabdul}) 

\begin{equation}

\left| {{f_j}(x)} \right| \le \mathop{\rm const} \cdot \bigl( {x + {( - 1)}^j q} \bigr)^{ - \varepsilon }, \quad 0 \le \varepsilon  < 1, \quad j = 1, 2. \label{eq:26.j}

\end{equation}

Заметим, что интегральные уравнения \eqref{eq:22.1} и \eqref{eq:22.2} можно переписать в следующем виде:

$$

{\tau '_1}(x) +  \int _x^1 {{\tau '_1}(t){K_1}(x,t)dt}  = f_1(x),$$ 

$${\tau '_2}(x) +  \int _{ - 1}^x {{\tau '_2}(t){K_2}(x,t)dt}  = f_2(x),$$

где  

$$

{K_1}(x,t) = 

\left\{\begin{array}{cl}

{K_{11}}(x,t), & x \le t \le ({{x + 1}})/{2},\\

{K_{11}}(x,t) + {K_{12}}(x,t), & ({{x + 1}})/{2} \le t \le 1;

\end{array}\right. $$

$$

{K_2}(x,t) = \left\{\begin{array}{cl}

{K_{22}}(x,t), & ({{x - 1}})/{2} \le t \le x;\\

{K_{21}}(x,t) + {K_{22}}(x,t), & - 1 \le t \le ({{x - 1}})/{2}.

\end{array}\right.

$$



Разрешая интегральные уравнения \eqref{eq:22.1} и \eqref{eq:22.2} методом последовательных приближений, находим

\begin{multline}

{\tau '_1}(x) = \frac{{2{a_1}(x) - 1}}{{1 - 2{b_1}(x)}}\nu _1^ - (x) +  \int _x^1 {\frac{{2{a_1}(t) - 1}}{{1 - 2{b_1}(t)}}{\Re _{11}}(x,t)\nu _1^ - (t)dt} +

\\

+  \int _{({{x + 1}})/{2}}^1 \frac{{2{a_1}(t) - 1}}{{1 - 2{b_1}(t)}}{\Re _{12}}(x,t)\nu _1^ - (t)dt +  {F_{11}}(x), \label{eq:27.1}

\end{multline}

\begin{multline}

{\tau '_2}(x) = \frac{{2{a_2}(x) + 1}}{{1 - 2{b_2}(x)}}\nu _2^ - (x) +  \int _{ - 1}^x \frac{{2{a_2}(t) + 1}}{{1 - 2{b_2}(t)}}{\Re _{21}}(x,t)\nu _2^ - (x)dt+

\\

+ \int _{ - 1}^{({{x - 1}})/{2}} \frac{{2{a_2}(t) + 1}}{{1 - 2{b_2}(t)}}{\Re _{22}}(x,t)\nu _2^ - (x)dt +  {F_{12}}(x), \label{eq:27.2}

\end{multline}

где ${F_{11}}(x)$ и ${F_{12}}(x)$ "--- известные функции, которые зависят от заданных функций ${\varphi _j}(x,y),$  

${g_j}(x),$  ${a_j}(x),$ ${b_j}(x),$ ${c_j}(x),$ ${d_j}(x)$ $( j = 1, 2)$ и $p_k(x),$ $r_k (x),$ а 

${\Re _{1j}}(x,t)$ и ${\Re _{2j}}(x,t)$ "--- резольвенты ядер ${K_{1j}}(x,t)$  и ${K_{1j}}(x,t)$ соответственно,  

причем

\begin{equation}

| \Re _{\sigma j} (x,t) | \le {\rm const}, \quad  

-1\leq (-1)^{\sigma}t\leq (-1)^{\sigma}x\leq -q, \quad  \sigma= 1, 2;  \; j = 1, 2; \label{eq:28.j}

\end{equation}

\begin{equation}

| F_{1j}(x) | \le {\rm const},\quad  -1\leq(-1)^{j}x\leq -q, \quad  j = 1, 2. \label{eq:28.1.j}

\end{equation}



\smallskip



\hypertarget{abd:B}{}



\texttt{Случай B.}

Пусть ${b_j}(x) \equiv  - {1}/{2}$ и ${c_j}(x) \ne 0$ $(j = 1, 2),$ 

тогда из \eqref{eq:13.1} и~\eqref{eq:13.2} получим интегральные уравнения Вольтерры второго рода относительно~${\tau _j}(x)$:

\begin{multline}

{\tau _1}(x) +  \int _x^1 {{{\bar K}_{11}}(x,t){\tau _1}(t)dt}  +  \int _{({{x + 1}})/{2}}^1 {{{\bar K}_{12}}(x,t){\tau _1}(t)dt} =

\\

 = \frac{{1 - 2{a_1}(x)}}{{2{c_1}(x)}}\nu _1^ - (x) - \frac{{{d_1}(x)}}{{{c_1}(x)}}, \label{eq:29.1}

\end{multline}

\begin{multline}

{\tau _2}(x) +  \int _{ - 1}^x {{{\bar K}_{21}}(x,t){\tau _2}(t)dt}  +  \int _{ - 1}^{({{x - 1}})/{2}} {{{\bar K}_{22}}(x,t){\tau _2}(t)dt} =

\\  =  - \frac{{2{a_2}(x) + 1}}{{{c_2}(x)}}\nu _2^ - (x) - \frac{{{d_2}(x)}}{{{c_2}(x)}},\label{eq:29.2}

\end{multline}

где

$$

{\bar K_{11}}(x,t) =  \int _x^t {\sum\limits_{k = 1}^n {\frac{{{{(t - z)}^{{\alpha _k} - 1}}}}{{4{c_1}(x)\Gamma ({\alpha _k})}}} {p_k}(z)dz},$$

$$ 

{\bar K_{12}}(x,t) =  -  \int _{({{x + 1}})/{2}}^t {\sum\limits_{k = 1}^n {\frac{{{{(t - z)}^{{\alpha _k} - 1}}}}{{4{c_1}(x)\Gamma ({\alpha _k})}}} {p_k}(z)dz},

$$

        $${\bar K_{21}}(x,y) =  -  \int _t^x {\sum\limits_{k = 1}^n {\frac{{{{(z - t)}^{{\beta _k} - 1}}}}{{4{c_2}(x)\Gamma ({\beta _k})}}} {r_k}(z)dz},$$

$$

{\bar K_{22}}(x,t) =  \int _t^{({{x - 1}})/{2}} {\sum\limits_{k = 1}^n {\frac{{{{(z - t)}^{{\beta _k} - 1}}}}{{4{c_2}(x)\Gamma ({\beta _k})}}} {r_k}(z)dz},$$

причём 

$$| {{{\bar K}_{\sigma j}}(x,t)} | \le {\rm const}, \quad 

-1\leq (-1)^{\sigma}t\leq (-1)^{\sigma}x\leq -q, \quad \sigma=1, 2.$$

Разрешая интегральные уравнения \eqref{eq:29.1} и \eqref{eq:29.2}, находим ${\tau _j}(x)$:

\begin{multline}

{\tau _1}(x) = \frac{{1 - 2{a_1}(x)}}{{2{c_1}(x)}}\nu _1^ - (x) +  \int _x^1 \frac{{1 - 2{a_1}(t)}}{{2{c_1}(t)}}{{\bar \Re }_{11}}(x,t)\nu _1^ - (t)dt +

\\

+  \int _{({{x + 1}})/{2}}^1 \frac{{1 - 2{a_1}(t)}}{{2{c_1}(t)}}{{\bar \Re }_{12}}(x,t)\nu _1^ - (t)dt +  {F_{21}}(x), \label{eq:30.1}

\end{multline}

\begin{multline}

{\tau _2}(x) =  - \frac{{2{a_2}(x) + 1}}{{{c_2}(x)}}\nu _2^ - (x) +  \int _{ - 1}^x \frac{{2{a_2}(t) + 1}}{{{c_2}(t)}}{{\bar \Re }_{21}}(x,t)\nu _2^ - (x)dt +

\\

+  \int _{ - 1}^{({{x - 1}})/{2}} \frac{{2{a_2}(t) + 1}}{{{c_2}(t)}}{{\bar \Re }_{22}}(x,t)\nu _2^ - (x)dt +  {F_{22}}(x), \label{eq:30.2}

\end{multline}

где ${F_{21}}(x)$ и ${F_{22}}(x)$ "--- известные функции, которые зависят от заданных функций, 

а~${\bar \Re _{1j}}(x,t)$ и ${\bar \Re _{2j}}(x,t)$ "--- резольвенты ядер ${\bar K_{1j}}(x,t)$ и ${\bar K_{2j}}(x,t)$ соответственно,  причем

\begin{equation}

| \bar \Re _{1j}(x,t) | \le {\rm const}, \quad 

| \bar \Re _{2j}(x,t) | \le {\rm const}, \quad   

| F_{2j}(x) | \le {\rm const}, \quad  j = 1, 2.\label{eq:31.j}

\end{equation}



\smallskip



\hypertarget{abd:C}{}



\texttt{Случай C.}

Пусть ${b_j}(x) \equiv  - {1}/{2}$  и ${c_j}(x) \equiv 0$ $(j = 1, 2),$ 

тогда из \eqref{eq:13.1} и \eqref{eq:13.2} получим

\begin{multline}

\nu _1^ - (x) =  \int _{({{x + 1}})/{2}}^1 {{\tau _1}(t)dt \int _{({{x + 1}})/{2}}^t {\sum\limits_{k = 1}^n {\frac{{{{(t - z)}^{{\alpha _k} - 1}}}}{{2\Gamma ({\alpha _k})\left( {2{a_1}(x) - 1} \right)}}} {p_k}(z)dz} }  - 

\\

-  \int _x^1 {{\tau _1}(t)dt \int _x^t {\sum\limits_{k = 1}^n {\frac{{{{(t - z)}^{{\alpha _k} - 1}}}}{{\Gamma ({\alpha _k})\left( {2{a_1}(x) - 1} \right)}}} {p_k}(z)dz} }  - \frac{{2{d_1}(x)}}{{2{a_1}(x) - 1}}, 

\label{eq:32.1}

\end{multline}

\begin{multline}

\nu _2^ - (x) =  \int _{ - 1}^x {{\tau _2}(t)dt \int _t^x {\sum\limits_{k = 1}^n {\frac{{{{(z - t)}^{{\beta _k} - 1}}}}{{2(2{a_2}(x) + 1)\Gamma ({\beta _k})}}{r_k}(z)dz} } }  - 

\\

-  \int _{ - 1}^{({{x - 1}})/{2}} {{\tau _2}(t)dt \int _t^{({{x - 1}})/{2}} {\sum\limits_{k = 1}^n {\frac{{{{(z - t)}^{{\beta _k} - 1}}}}{{2(2{a_2}(x) + 1)\Gamma ({\beta _k})}}{r_k}(z)dz} } }  - \frac{{2{d_2}(x)}}{{2{a_2}(x) + 1}}. \label{eq:32.2}

\end{multline}



Заметим, что решение задачи  Неймана для уравнения \eqref{eq:1.1} в области ${\Omega _0}$ с~краевыми условиями \eqref{eq:4.j} и

$${u_y}(x, +0) = {\nu _{1} ^+} (x), \quad (x, 0) \in {A_1}{B_1},$$

$${u_y}(x, +0) = {\nu _{2} ^+} (x),\quad (x, 0) \in {A_2}{B_2}$$

единственно и представимо в виде \cite{Islabd:bib:Islomabdul}:

\begin{multline}

u(x,y) =  \int _{{\sigma _1}} {{\varphi _1}(\xi ,\eta )\frac{\partial }{{\partial n}}G(\xi ,\eta ;x,y)dS}  -  \int _{{\sigma _2}} {{\varphi _2}(\xi ,\eta )\frac{\partial }{{\partial n}}G(\xi ,\eta ;x,y)dS}  + 

\\

 -  \int _q^1 {\nu _1^ + (t)G(t, 0; x, y)dt}  +  \int _{ - 1}^{ - q} \nu _2^ + (t)G(t, 0;x,y)dt, 

 \label{eq:33}

\end{multline}

где 

\begin{equation}

G(\xi ,\eta ;x,y) = \frac{1}{{2\pi }}\ln 

\biggl| {\frac{{{\theta _1}\bigl( {\frac{{\ln v + \ln \overline \mu  }}{{2\pi ir}}} \bigr){\theta _1}\bigl( {\frac{{\ln \overline v  + \ln \overline \mu  }}{{2\pi ir}}} \bigr)}}{{{\theta _1}\bigl( {\frac{{\ln v - \ln \mu }}{{2\pi ir}}} \bigr){\theta _1}\bigl( {\frac{{\ln \overline v  - \ln \mu }}{{2\pi ir}}} \bigr)}}} \biggr| \label{eq:34}

\end{equation}

"--- функция Грина задачи Неймана для уравнения ${u_{xx}} + {u_{yy}} = 0$  в области ${\Omega _0}$~\cite{Islabd:bib:Islomabdul}.

Здесь  $v = \xi  + i\eta ,$ $ \overline v  = \xi  - i\eta ,$ $\mu  = x + iy,$ $\overline \mu   = x - iy,$

$r = \frac{1}{{\pi i}}\ln q$, ${i^2} =  - 1$, ${\theta _1}(\xi ) $ "--- тета"=функция.

Из  \eqref{eq:33} при $y = 0$ находим функциональные соотношения между  $\tau _1^ + (x)$ и $\nu _1^ + (x)$   на интервалах ${A_1}{B_1}$ и ${{A_2}{B_2}} $, принесенные из области ${\Omega _0}$:

\begin{multline}

\tau _1^ + (x) = {F_1}(x) -  \int _q^1 {\nu _1^ + (t)G(t, 0; x, 0)dt}  +  \\

+ \int _{ - 1}^{ - q} {\nu _2^ + (t)G(t, 0; x, 0)dt}, \quad   q \le x \le 1, \label{eq:35.1}

\end{multline}

\begin{multline}

\tau _2^ + (x) = {F_1}(x) -  \int _q^1 {\nu _1^ + (t)G(t, 0; x, 0)dt}  +  \\ 

+ \int _{ - 1}^{ - q} {\nu _2^ + (t)G(t, 0; x, 0)dt},\quad   - 1 \le x \le  - q. \label{eq:35.2} 

\end{multline}

Здесь

\begin{equation}

{F_j}(x) =  \int _{{\sigma _1}} {{\varphi _1}(\xi ,\eta )\frac{\partial }{{\partial n}}G(\xi ,\eta ;x, 0)dS}  -  \int _{{\sigma _2}} {{\varphi _2}(\xi ,\eta )\frac{\partial }{{\partial n}}G(\xi ,\eta ;x, 0)dS}. 

\label{eq:36}

\end{equation}



\smallskip





\Section{Исследование интегральных уравнений}

В случае \hyperlink{abd:A}{\tt А} дифференцированием \eqref{eq:35.1} и \eqref{eq:35.2} по $x$ получим

 \begin{equation}

 \tau'_j(x) =  \int _{ - 1}^{ - q} {\nu _2^ + (t)\frac{{\partial G(t, 0;x, 0)}}{{\partial x}}dt}  -  \int _q^1 {\nu _1^+ (t)\frac{{\partial G(t, 0;x, 0)}}{{\partial x}}dt}  + {F'_j}(x). \label{eq:37.j}

 \end{equation}

Исключая ${\tau '_j}(x)$ $(j = 1, 2)$ из соотношений \eqref{eq:27.1}, \eqref{eq:27.2} и \eqref{eq:37.j} с учетом условия сопряжения \eqref{eq:3.j}, получим систему интегральных уравнений

\begin{equation}

\hspace{-1mm}

\begin{array}{l} 

\displaystyle

\frac{{2{a_1}(x) - 1}}{{1 - 2{b_1}(x)}} \nu _1^ - (x) + \frac{1}{\pi } \int _q^1 {\nu _1^ - (t){K_1}(x,t)dt}  + \\

\displaystyle 

+ 

\int _x^1 \frac{{2{a_1}(t) - 1}}{{1 - 2{b_1}(t)}}{\Re _{11}}(x,t)\nu _1^ - (t)dt +    

\int _{({{x + 1}})/{2}}^1 {\frac{{2{a_1}(t) - 1}}{{1 - 2{b_1}(t)}}{\Re _{12}}(x,t)\nu _1^ - (t)dt} =

\\

\displaystyle 

\hspace{5cm}

  = \frac{1}{\pi } \int _{ - 1}^{ - q} {\nu _2^ - (t){K_2}(x,t)dt}  + {F'_1}(x) - {F_{11}}(x),

  \\[2mm]

  \displaystyle

\frac{{2{a_2}(x) + 1}}{{1 - 2{b_2}(x)}}\nu _2^ - (x) - \frac{1}{\pi } \int _{ - 1}^{ - q} {\nu _2^ - (t){K_2}(x,t)dt}  + \\

\displaystyle 

+  \int _{ - 1}^x {\frac{{2{a_2}(t) + 1}}{{1 - 2{b_2}(t)}}{\Re _{21}}(x,t)\nu _2^ - (t)dt}   +  \int _{ - 1}^{({{x - 1}})/{2}} {\frac{{2{a_2}(t) + 1}}{{1 - 2{b_2}(t)}}{\Re _{22}}(x,t)\nu _2^ - (t)dt}  = \\

\displaystyle 

\hspace{5cm}

=  - \frac{1}{\pi } \int _q^1 {\nu _1^ - (t){K_1}(x,t)dt}  + {F'_2}(x) - {F_{12}}(x),

\end{array} \label{eq:37}

\hspace{-3mm}

\end{equation}

где

\begin{multline}

{K_j}(x,t) = \frac{{2\ln |t|}}{{x{\lambda _j}(t)\ln q}} + \frac{1}{{\pi {\lambda _j}(t)}}

\Bigl( {\frac{1}{{t - x}} - \frac{t}{{1 - tx}}} \Bigr) + 

\\

 + \frac{1}{{\pi {\lambda _j}(t)}}\sum\limits_{n = 1}^\infty  

 \Bigl( {\frac{{{q^{2n}}}}{{t - {q^{2n}}x}} - \frac{{{q^{2n}}t}}{{1 - {q^{2n}}tx}} - \frac{{{q^{ - 2n}}t}}{{1 - {q^{ - 2n}}tx}} + \frac{{{q^{ - 2n}}}}{{t - {q^{2n}}x}}} \Bigr). \label{eq:38}

\end{multline}

Учитывая \eqref{eq:14.j}, систему \eqref{eq:37} перепишем в следующем виде:

\begin{equation}

\begin{array}{l}

\displaystyle

\nu _1^ - (x) +  \int _q^1 {{K_{31}}(x,t)\nu _1^ - (t)dt}  = \frac{1}{\pi } \int _{ - 1}^{ - q} {\nu _2^ - (t){K_2}(x,t)dt}  + F_1^ * (x),\\

\displaystyle

\nu _2^ - (t) +  \int _{ - 1}^{ - q} {{K_{32}}(x,t)\nu _2^ - (t)dt}  =  - \frac{1}{\pi } \int _q^1 {\nu _1^ - (t){K_1}(x,t)dt}  + F_2^*(x),

\end{array} 

\label{eq:39}

\end{equation}

где  

$$F_1^ * (x) = \frac{{1 - 2{b_1}(x)}}{{2{a_1}(x) - 1}}\left( {{F'_1}(x) - {F_{11}}(x)} \right),$$  

$$F_2^ * (x) = \frac{{1 - 2{b_2}(x)}}{{2{a_2}(x) + 1}}\left( {{F'_2}(x) - {F_{12}}(x)} \right);$$

$${K_{31}}(x,t) = 

\left\{ \begin{array}{cc}

\frac{{1 - 2{b_1}(x)}}{{2{a_1}(x) - 1}}

\Bigl( \frac{{2{a_1}(t) - 1}}{{1 - 2{b_1}(t)}} \bigl( {{\Re _{11}}(x,t) + {\Re _{12}}(x,t)} \bigr) +

\\ 

\hspace{4.5cm} + \frac{1}{\pi }{K_1}(x,t) \Bigr), &({{x + 1}})/{2} \le t \le 1;

\\[2mm]

\frac{{1 - 2{b_1}(x)}}{{2{a_1}(x) - 1}}\Bigl( {\frac{{2{a_1}(t) - 1}}{{1 - 2{b_1}(t)}} {\Re _{11}}(x,t) + \frac{1}{\pi }{K_1}(x,t)} \Bigr), & x \le t \le ({{x + 1}})/{2};\\[2mm]

\frac{1}{\pi }  \frac{{1 - 2{b_1}(x)}}{{2{a_1}(x) - 1}} {K_1}(x,t), & q \le t \le x;

\end{array} \right.$$

$$

{K_{32}}(x,t) = \left\{ 

\begin{array}{cc}

\frac{{1 - 2{b_2}(x)}}{{2{a_2}(x) + 1}}\Bigl( \frac{{2{a_2}(t) + 1}}{{1 - 2{b_2}(t)}} \bigl( {{\Re _{21}}(x,t) + {\Re _{22}}(x,t)} \bigr)-

\\

\hspace{4.5cm} 

 - \frac{1}{\pi }{K_2}(x,t) \Bigr), & - 1 \le t \le ({{x - 1}})/{2};\\[2mm]

\frac{{1 - 2{b_2}(x)}}{{2{a_2}(x) + 1}} \Bigl( \frac{{2{a_1}(t) + 1}}{{1 - 2{b_2}(t)}} {\Re _{21}}(x,t) - {K_2}(x,t) \Bigr), & ({{x - 1}})/{2} \le t \le x;\\

\frac{{ - 1}}{\pi } \frac{{1 - 2{b_2}(x)}}{{2{a_1}(x) + 1}} \cdot {K_2}(x,t), & x \le t \le  - q. 

\end{array} 

\right.

$$ 

Заметим, что каждое уравнение системы \eqref{eq:39} является сингулярным интегральным уравнением нормального типа, индекс которого равен нулю. 

Точно так же, как и в  работах \cite{Abd:bib:Abdullaev, Islabd:bib:Islomabdul}, 

эта система известным методом Карлемана"--~Векуа \cite{Musxel:bib:Musxelesh} сводится к системе интегральных уравнений Фредгольма второго рода, разрешимость которой следует из единственности решения задачи~\hyperlink{abd:taskI}{I}.



Разрешая первое уравнение системы \eqref{eq:39}, находим

\begin{multline}

\nu _1^ - (x) = F_{1}^ * (x) + \frac{1}{\pi } \int _{ - 1}^{ - q} {\nu _2^ - (t){K_2}(x,t)dt}  + 

\\

 +  \int _q^1 {\left( {F_{1}^ * (t) + \frac{1}{\pi } \int _{ - 1}^{ - q} {\nu _2^ - (z){K_2}(t,z)dz} } \right){\Re_{31}}(x,t)dt},

\end{multline}

т.е.

\begin{equation}

\nu _1^- (x) = F_{1}^*(x) +  \int _q^1 {F_1^*(t){\Re_{31}}(x,t)dt} + \frac{1}{\pi } \int _{ - 1}^{ - q} {\nu _2^ - (z)K_2^ * (x,z)dz}, \label{eq:40}

\end{equation}

где  ${\Re_{31}}(x,t)$ "--- резольвента ядра  ${K_{31}}(x,t),$ а

$$K_2^*(x,t) = {K_2}(x,t) +  \int _q^1 {{K_2}(x,t)} {R_{31}}(x,t)dt,$$

причем  

$$

| {{\Re_{31}}(x,t)} | \le {\rm const},  \quad  

| {K_2^*(x,t)} | \le \rm const.$$

Подставляя найденные значение $ \nu _1^{}(x)$, т.~е. \eqref{eq:40} во второе уравнение системы \eqref{eq:39}, получим интегральное уравнение относительно $\nu _2^ - (x)$:

\begin{equation}

\nu _2^ - (x) +  \int _{ - 1}^{ - q} {\nu _2^ - (t)K_2^ * (x,t)dt}  = \bar F_2^ * (x), \label{eq:41}

\end{equation}

где

\begin{gather}

\bar F_2^*(x) = F_2^*(x) - \frac{1}{\pi } \int _q^1 {F_{1}^*(t){K_1}(x,t)dt - \frac{1}{\pi } \int _q^1 {{K_1}(x,t)dt}  \int _q^1 {F_1^*(z){\Re_{31}}(t,z)dz} }, \\

K_1^ * (x,z) = K_{32}^{}(x,z) + \frac{1}{{{\pi ^2}}} \int _q^1 {{K_1}(x,t)K_2^*(t,z)dt}, 

\end{gather}

причем   

$$| {\bar F_2^ * (x)} | \le {\rm const}, \quad | {K_1^ * (x,t)} | \le \rm const.$$



После того как найдены ${\nu^{-} _j}(x)$, из \eqref{eq:27.1}, \eqref{eq:27.2} и \eqref{eq:3.j} находим  ${\tau _j}(x)$ и ${\nu^{+} _j}(x)$. 

Следовательно, решение задачи~\hyperlink{abd:taskI}{I} можно восстановить в области ${\Omega _0}$ как решение задачи Неймана 

 \eqref{eq:33},  в областях $\Delta _j$ $(j = 1, 2)$ "--- как решение задачи Коши"--~Гурса, а в областях $D_j$ $(j = 1, 2, 3)$ "--- как решение задачи Гурса \cite{Islabd:bib:Islomabdul}. 

Итак, для случая   \hyperlink{abd:A}{\tt А} теорема~\hyperlink{abd:th2}{2} доказана.



\smallskip



В случае \hyperlink{abd:B}{\tt B} исключим ${\tau _j}(x)$ из \eqref{eq:30.1}, \eqref{eq:30.2}, \eqref{eq:35.1} и \eqref{eq:35.2} с~учетом условия сопряжения:

\begin{equation}

\begin{array}{l}

\displaystyle

\frac{{1 - 2{a_1}(x)}}{{2{c_1}(x)}}\nu _1^ - (x) + \frac{1}{\pi } \int _q^1 {\nu _1^ - (t){{\tilde K}_1}(x,t)dt}  + 

\\

\displaystyle

+  \int _x^1 \frac{{1 - 2{a_1}(t)}}{{2{c_1}(t)}}{{\bar \Re }_{11}}(x,t)\nu _1^ - (t)dt +   \int _{({{x + 1}})/{2}}^1 {\frac{{1 - 2{a_1}(t)}}{{2{c_1}(t)}}{{\bar \Re }_{12}}(x,t)\nu _1^ - (t)dt} =

\\

\displaystyle \hspace{4.5cm}

  = \frac{1}{\pi } \int _{ - 1}^{ - q} {\nu _2^ - (t){{\tilde K}_2}(x,t)dt}  + {F_1}(x) - {F_{21}}(x),\\[2mm]

\displaystyle 

\frac{{2{a_2}(x) + 1}}{{{c_2}(x)}}\nu _2^ - (x) + \frac{1}{\pi } \int _{ - 1}^{ - q} {\nu _2^ - (t){{\tilde K}_2}(x,t)dt}  - 

\\

\displaystyle

-

 \int _{ - 1}^x \frac{{2{a_2}(t) + 1}}{{{c_2}(t)}}{{\bar \Re }_{21}}(x,t)\nu _2^ - (x)dt  -  \int _{ - 1}^{({{x - 1}})/{2}} {\frac{{2{a_2}(t) + 1}}{{{c_2}(t)}}{{\bar \Re }_{22}}(x,t)\nu _2^ - (x)dt} = 

\\ 

\displaystyle \hspace{4.5cm}

   = \frac{1}{\pi } \int _q^1 {\nu _1^ - (t){{\tilde K}_1}(x,t)dt}  + {F_{22}}(x) - {F_2}(x).

\end{array}

\label{eq:42}

\end{equation}



\newpage

\noindent

Здесь

\begin{multline}

{\tilde K_j}(x,t) = \frac{{G(x, 0;t, 0)}}{{{\lambda _j}(t)}} = \frac{{\ln x\ln t}}{{{\lambda _j}(t)\ln q}} +

\\

 + \frac{1}{{\pi {\lambda _j}(t)}}\biggl( {\ln \Bigl| {\frac{{1 - tx}}{{t - x}}} \Bigr| + \sum\limits_{n = 1}^\infty  {\ln \Bigl| {\frac{{1 - {q^{2n}}tx}}{{t - {q^{2n}}x}} \cdot \frac{{{q^{2n}} - tx}}{{{q^{2n}}t - x}}} \Bigr|} } \biggr).

\end{multline}



Как и в случае \hyperlink{abd:A}{\tt A}, учитывая, что 

$$2{a_j}(x) + {( - 1)^j} > 0,$$ систему \eqref{eq:42}  перепишем в виде

\begin{equation}

\begin{array}{l}

\displaystyle

\nu _1^ - (x) +  \int _q^1 {{K_{41}}(x,t)\nu _1^ - (t)dt} =\\

\displaystyle

\hspace{3cm}

  = \frac{1}{\pi }\frac{{2{c_1}(x)}}{{1 - 2{a_1}(x)}} \int _{ - 1}^{ - q} {\nu _2^ - (t){{\tilde K}_2}(x,t)dt}  + F_3^ * (x),\\

\displaystyle

\nu _2^ - (t) +  \int _{ - 1}^{ - q} {{K_{42}}(x,t)\nu _2^ - (t)dt} =

\\

\displaystyle

\hspace{3cm}

  = \frac{1}{\pi }\frac{{2{c_2}(x)}}{{1 + 2{a_2}(x)}} \int _q^1 {\nu _1^ - (t){{\tilde K}_1}(x,t)dt}  + F_4^*(x),

\end{array}   \label{eq:43}

\end{equation}

где 

$$F_3^ * (x) = \frac{{2{c_1}(x)}}{{1 - 2{a_1}(x)}}\bigl( {{F_1}(x) - {F_{11}}(x)} \bigr),$$ 

$$F_4^ * (x) = \frac{{2{c_2}(x)}}{{2{a_2}(x) + 1}}\bigl( {{F_{22}}(x) - {F_2}(x)} \bigr);$$

$${K_{41}}(x,t) = \left\{ 

\begin{array}{cc}

\frac{{2{c_1}(x)}}{{1 - 2{a_1}(x)}}

\Bigl( \frac{{1 - 2{a_1}(t)}}{{2{c_1}(t)}} 

\bigl( {{\bar \Re }_{11}}(x,t) + {{\bar \Re }_{12}}(x,t) \bigr) +

\\

\hspace{4.5cm} 

 + \frac{1}{\pi }{{\tilde K}_1}(x,t) \Bigr), &({{x + 1}})/{2} \le t \le 1;

 \\[2mm]

\frac{{2{c_1}(x)}}{{1 - 2{a_1}(x)}}\Bigl( \frac{{1 - 2{a_1}(t)}}{{2{c_1}(t)}} {\Re _{11}}(x,t) + \frac{1}{\pi }{{\tilde K}_1}(x,t) \Bigr), & x \le t \le ({{x + 1}})/{2};\\

\frac{1}{\pi } \frac{{2{c_1}(x)}}{{1 - 2{a_1}(x)}}  {{\tilde K}_1}(x,t), & q \le t \le x;

\end{array} \right.$$

$$

{K_{42}}(x,t) = \left\{ 

\begin{array}{lc}

\frac{{2{c_2}(x)}}{{1 + 2{a_2}(x)}}

\Bigl( 

\frac{1}{\pi }{{\tilde K}_2}(x,t) - & \\

\hspace{1cm}

- \frac{{1 + 2{a_2}(t)}}{{2{c_2}(t)}} 

\bigl( {{{\bar \Re }_{21}}(x,t) + {{\bar \Re }_{22}}(x,t)} \bigr) \Bigr), & - 1 \le t \le ({{x - 1}})/{2};

\\[2mm]

\frac{{2{c_2}(x)}}{{1 + 2{a_2}(x)}}\Bigl( {\frac{1}{\pi }{{\tilde K}_2}(x,t) - \frac{{1 + 2{a_2}(t)}}{{2{c_2}(t)}} {{\bar \Re }_{21}}(x,t)} \Bigr), & ({{x - 1}})/{2} \le t \le x;\\[2mm]

\hspace{2cm}

\frac{1}{\pi } \frac{{2{c_2}(x)}}{{1 + 2{a_2}(x)}} {{\tilde K}_1}(x,t), & x \le t \le  - q.

\end{array} \right.

$$



Очевидно, что система \eqref{eq:43}  является системой интегральных уравнений Фредгольма второго рода, 

разрешимость которой следует из единственности решения задачи~\hyperlink{abd:taskI}{I}.

Теорема~\hyperlink{abd:th2}{2} далее доказывается, как в случае~\hyperlink{abd:A}{\tt А}.

  



\smallskip  

  



В случае~\hyperlink{abd:C}{\tt C} подставим \eqref{eq:13.1}, \eqref{eq:13.2} в \eqref{eq:35.1} и \eqref{eq:35.2}  соответственно и с~учетом условия \eqref{eq:3.j} получим

\begin{multline}

\tau _j (x) =  \int _q^1 {\tau _1(z)dz}  \int _q^z {{\lambda _1}(t)G(t, 0; x, 0)dt}  \int _t^z {\sum\limits_{k = 1}^n {\frac{{{{(z - s)}^{{\alpha _k} - 1}}}}{{2\Gamma ({\alpha _k})\left( {2{a_1}(z) - 1} \right)}}} {p_k}(s)ds}  - 

\\

 -  \int _{({{q + 1}})/{2}}^1 {\tau _1(z)dz}  \int _q^{2z - 1} {\lambda _1}(t)G(t, 0; x, 0 )dt 

\times \hspace{3cm}

 \\

 \hspace{3cm}

\times 

  \int _{({{t + 1}})/{2}}^z {\sum\limits_{k = 1}^n {\frac{{{{(z - s)}^{{\alpha _k} - 1}}}}{{2\Gamma ({\alpha _k})\left( {2{a_1}(z) - 1} \right)}}} {p_k}(s)ds}  + 

 \\

 +  \int _{ - 1}^{ - q} {\tau _2(z)dz}  \int _z^{ - q} {\lambda _2}(t)G(t, 0; x, 0)dt  \int _z^t {\sum\limits_{k = 1}^n {\frac{{{{(s - z)}^{{\beta _k} - 1}}}}{{2(2{a_2}(t) + 1)\Gamma ({\beta _k})}}{r_k}(s)} ds}  - 

\\

-  \int _{ - 1}^{ - ({{q + 1}})/{2}} {\tau _2(z)dz}  \int _{2z + 1}^{ - q} {{\lambda _2}(t)G(t, 0; x, 0)dt}  \times 

 \hspace{2cm}

\\

\times

 \int _z^{({{t - 1}})/{2}} {\sum\limits_{k = 1}^n {\frac{{{{(s - z)}^{{\beta _k} - 1}}}}{{2(2{a_2}(t) + 1)\Gamma ({\beta _k})}}{r_k}(s)} ds}  + {\tilde F_j}(x).

\end{multline}

Далее после некоторых упрощений получим систему интегральных уравнений Фредгольма второго рода относительно $\tau _j(x)$:

\begin{equation}

\begin{array}{ll}

\displaystyle

\tau _1(x) =  \int _q^1 {\tau _1(z){K_{51}}(x,z)dz}  +  \\ 

\displaystyle \hspace{4cm}

+ \int _{ - 1}^{ - q} {\tau _2(z){K_{52}}(x,z)dz}  + {{\tilde F}_1}(x), & q \le x \le 1,\\

\displaystyle

\tau _2(x) =  \int _{ - 1}^{ - q} {\tau _2^{}(z){K_{52}}(x,z)dz}  + \\ 

\displaystyle \hspace{4cm}

+  \int _q^1 {\tau _1(z){K_{51}}(x,z)dz}  + {{\tilde F}_2}(x), & - 1 \le x \le  - q. \end{array}

\hspace{-3mm}  \label{eq:44}

\end{equation}

Здесь

$$

{K_{51}}(x,z) = \left\{ 

\begin{array}{l}

\sum\limits_{k = 1}^n {\frac{1}{{2\Gamma ({\alpha _k})}}}  {\displaystyle \int _q^z } {\frac{{G(t, 0; x, 0)}}{{2{a_1}(t) - 1}}{\lambda _1}(t)dt}  {\displaystyle \int _t^z } {{{(z - s)}^{{\alpha _k} - 1}}{p_k}(s)ds}  - 

\\

 - \sum\limits_{k = 1}^n {\frac{1}{{2\Gamma ({\beta _k})}}}  

{\displaystyle 

 \int _q^{2z - 1} } {\frac{{G(t, 0; x, 0)}}{{2{a_1}(t) - 1}}{\lambda _1}(t)dt} {\displaystyle \int _{({{t + 1}})/{2}}^z } {{{(z - s)}^{{\alpha _k} - 1}}{p_k}(s)ds} ,\\

 \hspace{8cm} ({{1 + q}})/{2} \le t \le 1;

 \\

\sum\limits_{k = 1}^n {\frac{1}{{2\Gamma ({\alpha _k})}}}  {\displaystyle \int _q^z } {\frac{{G(t, 0; x, 0)}}{{2{a_1}(t) - 1}}{\lambda _1}(t)dt} {\displaystyle \int _t^z } {{{(z - s)}^{{\alpha _k} - 1}}{p_k}(s)ds} ,\\

\hspace{8cm}  q  \le t \le ({{1 + q}})/{2};

\end{array} \right.

$$

$$

{K_{52}}(x,z) = \left\{ \begin{array}{l}

\sum\limits_{k = 1}^n {\frac{1}{{2\Gamma ({\beta _k})}}} 

{\displaystyle  \int _z^{ - q} } {\frac{{G(t, 0; x, 0)}}{{2{a_2}(t) + 1}}{\lambda _2}(t)dt} {\displaystyle  \int _z^t } {{{(s - z)}^{{\beta _k} - 1}}{r_k}(s)ds}  - \\

 - \sum\limits_{k = 1}^n {\frac{1}{{2\Gamma ({\beta _k})}}}  {\displaystyle  \int _{2z + 1}^{ - q} } {\frac{{G(t, 0; x, 0)}}{{2{a_2}(t) + 1}}{\lambda _2}(t)dt} {\displaystyle  \int _z^{({{t - 1}})/{2}} } {{{(s - z)}^{{\beta _k} - 1}}{r_k}(s)ds} ,\\

\hspace{7.4cm}  - 1 \le t \le  - ({{1 + q}})/{2};

\\

{\displaystyle   \int _z^{ - q} } {{\lambda _2}(t)G(t, 0; x, 0)dt}  {\displaystyle   \int _z^t } {\sum\limits_{k = 1}^n {\frac{{{{(s - z)}^{{\beta _k} - 1}}}}{{2(2{a_2}(t) + 1)\Gamma ({\beta _k})}}{r_k}(s)} ds} ,\\ 

\hspace{7.4cm} - ({{1 + q}})/{2} \le t \le  - q;

\end{array} 

\right.

$$

\begin{multline}

{\tilde F_1}(x) = {F_1}(x) +  \int _q^1 \frac{{2{d_1}(t)}}{{2{a_1}(t) - 1}}G(t, 0; x, 0)dt -  \\ 

-  \int _{ - 1}^{ - q} \frac{{2{d_2}(t)}}{{2{a_2}(t) + 1}}G(t, 0; x, 0)dt, \quad   q \le x \le 1,

\end{multline}

\begin{multline}

{\tilde F_2}(x) = {F_2}(x) +  \int _q^1 \frac{{2{d_1}(t)}}{{2{a_1}(t) - 1}}G(t, 0; x, 0)dt - \\ 

-   \int _{ - 1}^{ - q} {\frac{{2{d_2}(t)}}{{2{a_2}(t) + 1}}G(t, 0; x, 0)dt}, \quad   - 1 \le x \le  - q. 

\end{multline}

причем  

$$| {{K_{5j}}(x,z)} | \le {\rm const},\quad  | {{{\tilde F}_j}(x)} | \le {\rm const},\quad  (j=1, 2).$$



Разрешая систему \eqref{eq:44}, однозначно определим $\tau _j(x)$, следовательно, из функциональных соотношений  \eqref{eq:13.1}, \eqref{eq:13.2} c учетом \eqref{eq:3.j} находим $\nu^{-} _j(x)$ и $\nu^{+} _j(x)$. 

Далее, как и в случаях~\hyperlink{abd:A}{\tt А} и~\hyperlink{abd:B}{\tt В},  однозначно восстанавливается  решение задачи~\hyperlink{abd:taskI}{I} в~каждой подобласти.



Таким образом, теорема~\hyperlink{abd:th2}{2} полностью доказана.



\smallskip



\thanks{{\bf ORCID}



{\bf  Обиджон Хайруллаевич Абдуллаев:} \url{http://orcid.org/0000-0001-8503-1268}



}





\RusRef



\begin{thebibliography}{99}







\RBibitem {Naxush:bib:Naxushev}

\zmath{https://zbmath.org/?q=an:0349.45011}

\mathscinet{http://www.ams.org/mathscinet-getitem?mr=0425554}

\by Нахушев~А.~М. 

\paper О~задаче Дарбу для одного вырождающегося нагруженного интегро-дифференциального  уравнения второго порядка

\jour Диффер. уравн.

\yr 1976

\vol 12

\issue 1

\pages 103--108

\mathnet{http://mi.mathnet.ru/de2654}







\RBibitem {Nax:bib:Naxushe}

\by Нахушев~А.~М. 

\paper Нагруженные уравнения и их приложения

\jour Диффер. уравн.

\yr  1983

\vol 19

\issue 1

\pages 86–94

\zmath{https://zbmath.org/?q=an:0536.35080}



\RBibitem{Naxusha:bib:Naxushevam}

\by Нахушев~А.~М.

\book Нагруженные уравнения и их применения 

\publaddr М.

\publ Наука 

\yr 2012 

\totalpages 233 







\RBibitem {Eleev:bib:Eleeva}

\zmath{https://zbmath.org/?q=an:0826.35079}

\by Елеев~В.~А.

\paper О некоторых краевых задачах для смешанно-нагруженных уравнений второго и~третьего порядка

\jour Диффер. уравн.

\vol 30

\issue 2

\yr 1994

\pages 230–236



\RBibitem{Elev:bib:Eleevdzrox}

\by Дзарахохов~А.~В., Елеев~В.~А.

\paper Об одной нелокальной краевой задаче для нагруженного уравнения третьего порядка

\jour Владикавк. матем. журн.

\yr 2004

\vol 6

\issue 3

\pages 36--46

\mathnet{http://mi.mathnet.ru/vmj212}

\mathscinet{http://www.ams.org/mathscinet-getitem?mr=2120840}

\zmath{https://zbmath.org/?q=an:1113.35330}







\RBibitem {Kaz:bib:Kaziyev}

\zmath{https://zbmath.org/?q=an:0371.45009|an:0395.45019}

\by Казиев~В.~М.

\paper О~задаче Дарбу для одного нагруженного интегро-дифференциального уравнения второго порядка

\jour Диффер. уравн.

\yr 1978

\vol 14

\issue 1

\pages 181--184

\mathnet{http://mi.mathnet.ru/de3294}



\RBibitem {Kaze:bib:Kaziev}

\by Казиев~В.~М.

\paper Задача Гурса для одного нагруженного интегро-дифференциального уравнения

\jour Диффер. уравн.

\yr 1981

\vol 17

\issue 2

\pages 313--319

\mathnet{http://mi.mathnet.ru/de4195}

\zmath{https://zbmath.org/?q=an:0487.45009}







\RBibitem {Lan:bib:Lanin}

\by Ланин~И.~Н.

\paper Краевая задача для одного нагруженного гиперболо-параболического уравнения третьего порядка

\jour Диффер. уравн.

\yr 1981

\vol 17

\issue 1

\pages 97--106

\mathnet{http://mi.mathnet.ru/de4173}

\zmath{https://zbmath.org/?q=an:0564.35077}



\RBibitem{Islom:bib:Islomov}

\zmath{https://zbmath.org/?q=an:0940.35144}

\paper Об одной краевой задаче для нагруженного уравнения второго порядка

\by Исломов~Б.~И., Курьязов~Д.~М.

\jour ДАН РУз

\yr  1996

\issue  1--2

\pages 3--6



\RBibitem {Kur:bib:Kuryaz}

\by Курьязов~Д.~М.

\paper Краевая задача для нагруженного уравнения смешанного типа с двумя линиями вырождения

\jour УзМЖ

\yr 1999

\issue 5

\pages 40--46









\RBibitem {Ram:bib:Ramaz}

\zmath{https://zbmath.org/?q=an:05267907}

\by Рамазанов~М.~И.

\paper О нелокальной задаче для нагруженного гиперболо-эллиптического уравнения в прямоугольной области

\elink \url{http://www.math.kz/images/journal/2002-4/Ramazanov.pdf}

\jour Математический журнал. Алматы

\vol 2

\issue 4

\pages 75--81

\yr 2002



\RBibitem {Xub:bib:Xubiyev}

\jour Доклады Адыгской (Черкесской) Международной академии наук

\vol 7

\issue 2 

\yr 2005

\pages 74--77

\paper Об одной краевой задаче для нагруженного уравнения смешанного гиперболо-параболического типа

\by Хубиев~К.~У.

\elib{http://elibrary.ru/item.asp?id=11152898}







\RBibitem {Sab:bib:Sabit}

\by Сабитов~К.~Б., Мелишева~Е.~П.

\paper Задача Дирихле для нагруженного уравнения смешанного типа в~прямоугольной области

\jour Изв. вузов. Матем.

\yr 2013

\issue 7

\pages 62--76

\mathnet{http://mi.mathnet.ru/ivm8811}

%\transl

%\jour Russian Math. (Iz. VUZ)

%\yr 2013

%\vol 57

%\issue 7

%\pages 53--65

%\crossref{10.3103/S1066369X13070062}

%\scopus{http://www.scopus.com/record/display.url?origin=inward&eid=2-s2.0-84879752682}



\RBibitem {Sab:bib:sabitov}

\by Сабитов~К.~Б.

\paper Начально-граничная задача для параболо-гиперболического уравнения с~нагруженными слагаемыми

\jour Изв. вузов. Матем.

\yr 2015

\issue 6

\pages 31--42

\mathnet{http://mi.mathnet.ru/ivm9008}

%\transl

%\jour Russian Math. (Iz. VUZ)

%\yr 2015

%\vol 59

%\issue 6

%\pages 23--33

%\crossref{10.3103/S1066369X15060055}

%\scopus{http://www.scopus.com/record/display.url?origin=inward&eid=2-s2.0-84938234834}





\RBibitem {Meli:bib:Melish}

\by Мелишева~Е.~П.

\paper Задача Дирихле для нагруженного уравнения Лаврентьева--Бицадзе

\jour Вестн. СамГУ. Естественнонаучн. сер.

\yr 2010

\issue 6(80)

\pages 39--47

\mathnet{http://mi.mathnet.ru/vsgu192}





%%%%%%







\Bibitem{Abdu:bib:Abdull}

\by Abdullayev~O.~Kh.

\paper About a method of research of the non-local problem for the loaded mixed type equation in double-connected domain

\jour Bulletin KRASEC. Phys. \& Math. Sci.

\yr 2014

\vol 9

\issue 2

\pages 3--12

\mathnet{http://mi.mathnet.ru/vkam47}

\crossref{10.18454/2313-0156-2014-9-2-3-12}



\RBibitem{Abd:bib:Abdullaev}

\by Абдуллаев~О.~Х.

\paper Краевая задача для нагруженного уравнения эллиптико-гиперболического типа в двусвязной области

\jour Вестник КРАУНЦ. Физ.-мат. науки

\yr 2014

\issue 1(8)

\pages 33--48

\mathnet{http://mi.mathnet.ru/vkam27}

\crossref{10.18454/2079-6641-2014-8-1-33-48}

\elib{http://elibrary.ru/item.asp?id=21624158}



 

 \RBibitem{Islabd:bib:Islomabdul}

\by Исломов~Б.~И., Абдуллаев~О.~Х.

\paper Краевая задача типа задачи Бицадзе для уравнения третьего порядка эллиптико-гиперболического типа в двусвязной области

\jour Доклады Адыгской (Черкесской) Международной академии наук

\yr 2004

\vol 7

\issue 1

\pages 42--46



\RBibitem {Bits:bib:Bitsad}

\zmath{https://zbmath.org/?q=an:03238931}

\by Бицадзе~А.~В.

\book Краевые задачи эллиптических уравнений второго порядка

\publaddr М.

\publ Наука

\yr 1966

\totalpages 203



















\RBibitem {Musxel:bib:Musxelesh}

\by Мусхелишвили~Н.~И.

\book Сингулярные интегральные уравнения. Граничные задачи теории функций и некоторые их приложения к~математической физике

\zmath{https://zbmath.org/?q=an:03277892}

\publaddr М.

\publ Наука

\yr 1968

\totalpages 513 









\end{thebibliography}



\RuDate



\smallskip



\makeentitle





\smallskip



\thanks{{\bf ORCID}



{\bf  Obidjon Kh. Abdullayev:} \url{http://orcid.org/0000-0001-8503-1268}



}





\EngRef



\begin{thebibliography}{99}







\Bibitem {Naxush:bib:Naxushev:e}

\lang In Russian

\by Nahushev~A.~M.

\paper The Darboux problem for a certain degenerate second order loaded integrodifferential equation

\jour Differ. Uravn.

\vol 12 

\yr 1976

\issue 1

\pages 103–108







\Bibitem {Nax:bib:Naxushe:e}

\zmath{https://zbmath.org/?q=an:0536.35080}

\by Nakhushev~A.~M.

\paper Loaded equations and their applications

\jour Differ. Equations 

\vol 19

\issue 1

\pages 74--81 

\yr 1983



\Bibitem{Naxusha:bib:Naxushevam:e}

\lang In Russian

\by Nakhushev~A.~M.

\book Nagruzhennye uravneniia i ikh primeneniia  \rm [Loaded equations and their applications]

\publaddr Moscow

\publ Nauka 

\yr 2012 

\totalpages 233 



\Bibitem {Eleev:bib:Eleeva:e}

\zmath{https://zbmath.org/?q=an:0826.35079}

\by Eleev~V.~A.

\paper Some boundary value problems for mixed loaded equations of second and third order

\jour Differ. Equations 

\vol 30

\issue 2

\pages 210--217 

\yr 1994



\Bibitem{Elev:bib:Eleevdzrox:e}

\lang In Russian

\by Dzarakhohov~A.~V., Eleev~V.~A.

\paper On a nonlocal boundary value problem for a third-order loaded equation

\jour Vladikavkaz. Mat. Zh.

\yr 2004

\vol 6

\issue 3

\pages 36--46







\Bibitem {Kaz:bib:Kaziyev:e}

\zmath{https://zbmath.org/?q=an:0371.45009|an:0395.45019}

\by Kaziev~V.~M.

\paper Darboux's problem for a loaded second-order integrodifferential equation

\jour Differ. Equations 

\vol 14

\issue 1

\pages  130--133 

\yr 1978



\Bibitem {Kaze:bib:Kaziev:e}

\by Kaziev~V.~M.

\paper Goursat's problem for a loaded integrodifferential equation

\jour Differ. Equations 

\vol 17

\pages 216--220 

\yr 1981

\issue 2

\zmath{https://zbmath.org/?q=an:0487.45009}







\Bibitem {Lan:bib:Lanin:e}

\by Lanin~I.~N.

\paper Boundary-value problem for a loaded third-order hyperbolic-parabolic equation

\jour Differ. Equations 

\vol 17

\issue 1

\pages 66--72

\yr 1981

\zmath{https://zbmath.org/?q=an:0564.35077}



\Bibitem{Islom:bib:Islomov:e}

\lang In Russian

\by Islamov~B.~I., Kuryazov~D.~M.

\paper On a boundary value problem for loaded equation of the second order 

\jour Dokl. Akad. Nauk Resp. Uzb.

\yr  1996

\issue  1--2

\pages 3--6



\Bibitem {Kur:bib:Kuryaz}

\lang In Russian

\by Kuryazov~D.~M.

\paper A boundary value problem for a loaded equation of mixed type with two lines of degeneracy

\jour Uzbekskii Matematicheskii Zhurnal

\yr 1999

\issue 5

\pages 40--46









\Bibitem {Ram:bib:Ramaz:e}

\lang In Russian

\by Ramazanov~M.~I.

\paper On a nonlocal problem for a loaded hyperbolic-elliptic equation in a rectangular domain

\jour Matematicheskii zhurnal. Almaty

\elink \url{http://www.math.kz/images/journal/2002-4/Ramazanov.pdf}

\vol 2

\issue 4

\pages 75--81

\yr 2002





\Bibitem {Xub:bib:Xubiyev:e}

\lang In Russian

\jour Doklady Adygskoi (Cherkesskoi) Mezhdunarodnoi akademii nauk

\vol 7

\issue 2 

\yr 2005

\pages 74--77

\paper On one boundary value problem for a loaded  mixed hyperbolic-parabolic type equation 

\by Khubiev~K.~U.









\Bibitem {Sab:bib:Sabit:e}

\paper The Dirichlet problem for a loaded mixed-type equation in a rectangular domain

\by Sabitov~K.~B., Melisheva~E.~P.

\jour Russian Math. (Iz. VUZ)

\yr 2013

\vol 57

\issue 7

\pages 53--65

\crossref{10.3103/S1066369X13070062}

\scopus{http://www.scopus.com/record/display.url?origin=inward&eid=2-s2.0-84879752682}



\Bibitem {Sab:bib:sabitov:e}

\paper Initial-boundary problem for parabolic-hyperbolic equation with loaded summands

\by Sabitov~K.~B.

\jour Russian Math. (Iz. VUZ)

\yr 2015

\vol 59

\issue 6

\pages 23--33

\crossref{10.3103/S1066369X15060055}

\scopus{http://www.scopus.com/record/display.url?origin=inward&eid=2-s2.0-84938234834}





\Bibitem {Meli:bib:Melish:e}

\lang In Russian

\by Melisheva~E.~P.

\paper Dirichlet problem for loaded equation of Lavrentiev--Bizadze

\jour Vestnik SamGU. Estestvenno-Nauchnaya Ser.

\yr 2010

\issue 6(80)

\pages 39--47







\Bibitem{Abdu:bib:Abdull:e}

\by Abdullayev~O.~Kh.

\paper About a method of research of the non-local problem for the loaded mixed type equation in double-connected domain

\jour Bulletin KRASEC. Phys. \& Math. Sci.

\yr 2014

\vol 9

\issue 2

\pages 3--12

\mathnet{http://mi.mathnet.ru/vkam47}

\crossref{10.18454/2313-0156-2014-9-2-3-12}



\Bibitem{Abd:bib:Abdullaev:e}

\lang In Russian

\by Abdullayev~O.~Kh.

\paper A boundary value problem for a loaded elliptic-hyperbolic type equation  in a double-connected domain

\jour Vestnik KRAUNTs. Fiz.-mat. nauki \rm [Bulletin KRASEC. Phys. \& Math. Sci.]

\yr 2014

\issue 1(8)

\pages 33--48

\crossref{10.18454/2079-6641-2014-8-1-33-48}





 

\Bibitem{Islabd:bib:Islomabdul:e}

\lang In Russian

\by Islomov~B.~I., Abdullayev~O.~Kh.

\paper A Boundary value problem of Bitsadze type problem for a third-order  elliptic-hyperbolic type equation in a double-connected domain

\jour Doklady Adygskoi (Cherkesskoi) Mezhdunarodnoi akademii nauk

\yr 2004

\vol 7

\issue 1

\pages 42--46

 



\Bibitem {Bits:bib:Bitsad:e}

\by Bitsadze~A.~V.

\book Boundary value problems for second order elliptic equations

\zmath{https://zbmath.org/?q=an:0167.09401}

\serial North-Holland Series in Applied Mathematics and Mechanics

\vol 5

\publaddr Amsterdam

\publ North-Holland Publ. 

\totalpages 211 

\yr 1968





\Bibitem {Musxel:bib:Musxelesh:e}

\by Muskhelishvili~N.~I.

\book Singular integral equations

\publaddr Groningen

\publ Wolters-Noordhoff Publ.

\totalpages  447 

\yr 1967

\zmath{https://zbmath.org/?q=an:03277891}









\end{thebibliography}



\EngDate



\Russian

\renewcommand{\baselinestretch}{0.9}



\newpage



%\end{document}

%

